\documentclass[11pt,a4paper]{article}
\usepackage[top=2.5cm,bottom=2.8cm,left=2.5cm,right=2.5cm]{geometry}
\usepackage{amsmath,amssymb,amsthm}
\usepackage{hyperref,xcolor,enumitem,booktabs,array}
\usepackage{graphicx,subcaption}

\hypersetup{colorlinks=true,linkcolor={blue!60!black},
  citecolor={blue!60!black},urlcolor={blue!60!black}}

\newtheorem{theorem}{Theorem}[section]
\newtheorem{lemma}[theorem]{Lemma}
\newtheorem{proposition}[theorem]{Proposition}
\newtheorem{corollary}[theorem]{Corollary}
\theoremstyle{definition}\newtheorem{definition}[theorem]{Definition}
\theoremstyle{remark}
\newtheorem{observation}[theorem]{Observation}
\newtheorem{remark}[theorem]{Remark}
\newtheorem{conjecture}[theorem]{Conjecture}

\newcommand{\vp}{\varphi}
\newcommand{\lam}{\lambda}
\newcommand{\Ja}{J_{2a}}
\newcommand{\Jfour}{J_4}
\newcommand{\TCMB}{T_{\mathrm{CMB}}}
\newcommand{\Lcond}{\Lambda_{\mathrm{cond}}}
\newcommand{\rCMB}{\rho_{\mathrm{CMB}}}
\newcommand{\SU}{\mathrm{SU}}
\newcommand{\Efb}{E_{\mathrm{fb}}}
\newcommand{\Kfb}{K_{\mathrm{fb}}}
\newcommand{\twoI}{2I}
\newcommand{\twoT}{2T}
\newcommand{\QH}{Q_H}
\newcommand{\vEW}{v_{\mathrm{EW}}}
\newcommand{\Ztwo}{\mathbb{Z}_2}
\newcommand{\Zthree}{\mathbb{Z}_3}
\newcommand{\dimq}[1]{d_{#1}}
\newcommand{\Tmat}{T}
\newcommand{\Jones}{V}
\newcommand{\lk}{\mathrm{lk}}
\newcommand{\phistar}{\phi^*}
\newcommand{\tang}{\dot\Gamma}

\title{\textbf{Dynamical Origin of Quark Masses in the\\
Density-Feedback Hopfion:\\
Analytic Constraints from the Trefoil Formation Event}
\\[0.6em]
{\large Companion Paper XIX to the Density-Feedback
Faddeev--Niemi Hopfion Series}
\\[0.6em]
\large Preprint --- July 2026}
\author{Frederick Manfredi}
\date{}

\begin{document}
\maketitle

\begin{abstract}
Paper~XVI~\cite{Paper16} documented a systematic negative result:
six candidate mechanisms for the quark generation mass hierarchy
(600-cell vertices, $\twoT$ 1-D irreps, $\SU(3)$ level, $\twoT/Q_8$
coset, generation-dependent colour shift, $E_8$ Coxeter labelling)
each fail for a specific structural reason.  The paper's own
Discussion identified two patterns: no mechanism alone
can supply generation-dependent input, and a dynamical correction
(beyond any representation-theoretic backbone) is required.
The present paper confirms and refines this picture.  Two facts
established in Paper~XVIII change what the calculation is: the
$\Zthree$ symmetry of the ideal $T_{2,3}$ trefoil is broken from the
outset by the single perturbation site of the formation event
(Remark~5.3 of Paper~XVIII~\cite{Paper18}), and %\ref{P18:rem:crossing_path}
the $\QH=3$ configuration is transient and unstable, so its
observables are properties of the formation-and-decay process
rather than of a static geometry.  Quark masses are therefore not
intrinsic energy scales of the trefoil but characteristic scales of
how a single incoming perturbation partitions across the trefoil
geometry during transit.
We derive four results within this framework.
(1)~\textbf{The minor radius $r_0=R_0/C_2^*$ is analytically fixed}
by the $\QH=2$ profile.  The Bogomolny profile
$f(\rho)=2\arctan(\rho^{-C_2^*})$ has slope $|f'(1)|=C_2^*$ at its
transition midpoint, giving wall width $\Delta\rho=1/C_2^*$; the
trefoil's minor radius equals this wall width by the satellite
construction.  This closes an open problem shared by Papers~XV
and~XVIII.
(2)~\textbf{The single-loss-coefficient cascade model is
analytically falsified.}  Neither the arc-segment density ratio
$\eta_B=\rho_B/\rho_A=0.81$ (Paper~XVI) nor the per-crossing
transmission $\eta_{\mathrm{cross}}=0.266$ reproduces observed
quark mass ratios within either isospin triplet.
(3)~\textbf{The uniform $\vp$-tower generation multiplier is
analytically falsified.} A hypothesis $\Delta m_g\propto\vp^{2g}$
predicts intra-triplet ratios $\vp^{4}\approx6.85$ that miss
observed ratios $(m_b/m_d\approx 890$, $m_t/m_u\approx 78600)$ by
two to four orders of magnitude.  The correct replacement structure,
identified from the empirical cross-ratio pairing of intra-triplet
generation gaps (Remark~\ref{P19:rem:cross_ratio}), is a geometric
ladder of four quark mass ratios with common ratio $\sqrt{\vp}$,
anchored at $m_s/m_d$
(Theorem~\ref{P19:thm:sqrt_phi_ladder}).  Combined with
Paper~XVI's derived $m_s=93.32\,\mathrm{MeV}$ ($0.09\%$) and the
framework-derived ratio $m_s/m_d=\sqrt{k/2}\cdot e^{8\pi/9}=19.991$
(equation~\eqref{P19:eq:spin_corrected_ratio}, within $0.047\%$ of
$n_e=20$), this yields
$m_b=4214\,\mathrm{MeV}$ ($0.7\%$ off PDG), $m_t/m_c$ at $6\%$,
and $m_c/m_u$ at $19\%$ (Corollary~\ref{P19:cor:absolute_masses}).
Notably, $m_s/m_d\approx n_e=20$, the lepton CMB tower level
derived independently from $T_{\mathrm{CMB}}$ and geometry;
the spin-Casimir identity $k=4C_2(\tfrac12)$, exact at $k=3$,
gives the $0.047\%$ residual from the integer (Open Problem~O6).
(4)~\textbf{The two-triplet mass-scale ratio is predicted at the
$20\%$ level} by the writhe-$\Ztwo$-weighted $z$-tangent
distribution: $M_{\mathrm{down}}/M_{\mathrm{up}}\approx0.189$ versus
observed $0.157$.  This is a parameter-free consequence of the
site-tangent geometry of Paper~XVI and the isospin identification
of Paper~XVIII (Conjecture~5.4 of that paper). %\ref{P18:conj:isospin}
The four results together sharpen the question: quark masses within
each triplet are set by the azimuthal-position-dependent
partition of the incoming perturbation energy, a calculation that
requires the dynamic integrator infrastructure rather than the
gradient-flow saddle.  A fifth result resolves Paper~XVIII's Open
Problem~O4: the multi-component Jones polynomial reaction rule is
not numerical equality of Jones products but unit-equivalence in
$\mathbb{Z}[\zeta_5]$; both production channels
($T(2,2)\sqcup T(2,2)\to T(2,3)\sqcup U$ and
$T(2,2)\sqcup U\to T(2,3)$) are verified algebraically, with
specific unit factors $\zeta_5+\zeta_5^4=1/\vp$ and
$-\vp\cdot\zeta_5$ respectively, reflecting the Hopf-link and
satellite Jones values $|\Jones_{T(2,2)}(q_5)|=|\Jones_{\rm sat}(q_5)|
=1/\vp$ established in Paper~XVIII's Theorem~2.3.
The reconciliation of Conjectures
\ref{P19:conj:six_quarks_r19} and~\ref{P19:conj:isospin_r19} into a single
dynamical picture is stated as a sharpened conjecture.
A cross-sector formal result (Proposition~\ref{P19:prop:solar_angle_jones})
establishes that the leading-order solar mixing angle satisfies
$\tan\theta_{12}^{(0)}=1/\vp$ via Paper~XVII's $\mathrm{SU}(2)_3$
S-matrix; the same modulus $1/\vp$ also equals
$|\Jones_{T(2,2)}(q_5)|$, a numerical coincidence recorded but not
established as a derivation (Remark~\ref{P19:rem:solar_physical}).
\end{abstract}

\tableofcontents
\bigskip\newpage

%══════════════════════════════════════════════════════════════════════════════
\section{Introduction}
\label{P19:sec:intro}
%══════════════════════════════════════════════════════════════════════════════

The quark mass hierarchy remains the framework's most conspicuous
unresolved problem.  Six orders of magnitude separate $m_u\approx
2.2\,\mathrm{MeV}$ from $m_t\approx 172\,\mathrm{GeV}$; four
orders separate $m_d\approx 4.7\,\mathrm{MeV}$ from $m_b\approx
4.2\,\mathrm{GeV}$.  No smooth function of a single generation
index reproduces both sequences, and Paper~XVI~\cite{Paper16}'s
systematic survey established that no combinatorial or
representation-theoretic mechanism internal to $\twoT$ can supply
the required generation-dependent input.  The strange quark alone
matches (to $0.09\%$, via a specific $E_8$ Coxeter exponent
assignment); the other five undershoot by factors of $1.7$ to $49$
in a systematic, always-positive-deficit pattern that Paper~XVI's
Discussion identified as the qualitative signature of a missing
$\QH=3$-specific, generation-dependent energy contribution.

The present paper reframes the problem in light of two facts
established in Paper~XVIII~\cite{Paper18}.

\paragraph{Fact 1: $\Zthree$ is broken from the outset.}  The
$\Zthree$ symmetry of the ideal trefoil is a construction fiction.
The $\QH=3$ trefoil is produced by threading a $\QH=1$ perturbation
through the neighbourhood $N(\gamma_1)$ of one component of the
$\QH=2$ Hopf link (Proposition~5.2 of %\ref{P18:prop:satellite}
Paper~XVIII~\cite{Paper18}); this threading event has a
definite azimuthal location, which fixes a distinguished direction
on the incipient trefoil.  The three crossings, three lobes, and
three midpoints that emerge in the resulting configuration are
labeled by their distance-along-the-tube from the threading site,
not by any residual $\Zthree$ symmetry.  A physical $\QH=3$
condensate carries a memory of its formation; only the average
over an ensemble of formations restores $\Zthree$ in expectation.

\paragraph{Fact 2: $\QH=3$ is transient.}  The satellite Jones
invariant of the trefoil at the WZW root $q_5=e^{2\pi i/5}$
satisfies $|\Jones_{\rm sat}(q_5)|=1/\vp$
(Paper~XVIII~\cite{Paper18}, Theorem~2.3). %\ref{P18:thm:trefoil_satellite_jones}
The formation path passes through an unfused intermediate
configuration with rational satellite invariant ($=1$), providing
no irrational BPS fixed point during formation; the energy
landscape has a flat plateau (Remark~10.2 of Paper~XVIII). %\ref{P18:rem:saddle_difficulty}
A physical $\QH=3$ condensate does not persist as an equilibrium
configuration; it forms, cascades, and decays over the collision
timescale.  Observable quantities associated with it are properties
of this transit, not of any equilibrium state.

\paragraph{Consequence.}  Quark masses in this framework are not
intrinsic energy scales of the trefoil geometry.  They are
characteristic scales of how a single incoming perturbation energy
partitions across the trefoil sites during its transit through the
formation-and-decay cascade.  This changes the calculation
qualitatively: instead of computing static site energies from the
condensate action and reading them as masses, one must compute the
dynamical partition of $E_{\mathrm{pert}}$ across the twelve
geometric sites, given the perturbation's azimuthal location and
the trefoil's Frenet geometry.

The present paper does not compute this partition directly.  That
calculation requires the dynamic integrator infrastructure
initiated in Papers~XVII and~XVIII (the suppression-weighted
restoring force applied to the $\QH=2$ vacuum), extended to
simulate the $\QH=2+\QH=1\to\QH=3$ formation event with the
velocity field $v(\mathbf{x},t)$ recorded at each of the twelve
sites.  What the present paper does is establish, analytically,
four constraints that any such calculation must satisfy or
falsify.

\paragraph{Organisation.}
Section~\ref{P19:sec:r0_derivation} closes the open problem shared by
Papers~XV and~XVIII on the physical origin of $r_0$: an analytic
derivation of $r_0=R_0/C_2^*$ from the $\QH=2$ profile function.
Section~\ref{P19:sec:reconciliation} reconciles Conjectures~5.1
and~5.4 of Paper~XVIII (six-quark combinatorics and writhe-based
isospin) into a single dynamical picture, resolving an internal
tension between them.
Section~\ref{P19:sec:cascade_falsification} formalises the
single-loss-coefficient cascade model and shows analytically that
it is falsified by observed quark mass ratios (Result~A of the
paper's analytic constraints).
Section~\ref{P19:sec:phi_tower_falsification} formalises the
$\vp^{2g}$ generation-multiplier hypothesis and shows analytically
that it is falsified (Result~B), then examines whether the surviving
$\bar n_{\mathrm{up}}/\bar n_{\mathrm{down}}\approx\vp$ structure has
a Jones-polynomial origin: the direct comparison does not reproduce
it, and the matter is left as an open question
(Proposition~\ref{P19:prop:jones_modulus_ratio}).
Section~\ref{P19:sec:tangent_ratio_prediction} derives the two-triplet
mass-scale ratio $M_{\mathrm{down}}/M_{\mathrm{up}}$ from the
writhe-$\Ztwo$-weighted $z$-tangent distribution and shows it
matches observation at the $20\%$ level (Result~C).
Section~\ref{P19:sec:multi_component_jones} develops the
multi-component Jones polynomial framework of Paper~XVIII's Open
Problem~O1, proves the correct reaction rule
(unit-equivalence in $\mathbb{Z}[\zeta_5]$, not numerical equality),
and verifies both $\QH=3$ production channels algebraically,
resolving Paper~XVIII's Open Problem~O4.
Section~\ref{P19:sec:experiment} collects experimental predictions
and constraints.
Section~\ref{P19:sec:open} states the sharpened open problems.
Section~\ref{P19:sec:summary} summarises.

%══════════════════════════════════════════════════════════════════════════════
\section{The minor radius $r_0 = R_0/C_2^*$ from the $\QH=2$ profile}
\label{P19:sec:r0_derivation}
%══════════════════════════════════════════════════════════════════════════════

The trefoil's minor tube radius $r_0=0.874$ was introduced in
Papers~XV and~XVI as a numerical construction parameter, without an
analytic derivation.  Paper~XV~\cite{Paper15}, Remark~14.5 %\ref{P15:rem:r0_origin}
identified the near-coincidence
\begin{equation}
  r_0 \;\approx\; \frac{R_0}{C_2^*} \;=\; \frac{3}{3.4318} \;=\; 0.874172,
  \label{P19:eq:r0_C2}
\end{equation}
matching the construction value to $0.020\%$, motivated by the
satellite construction (Paper~XVIII~\cite{Paper18},
Proposition~5.2): the trefoil is built inside $N(\gamma_1)$, %\ref{P18:prop:satellite}
so $r_0$ should inherit its scale from the $\QH=2$ tube geometry
rather than from a $\QH=3$-internal saddle condition.  Paper~XV
noted explicitly that this identification is
\emph{motivated} by the satellite construction but not
\emph{derived} from it: the map from $N(\gamma_1)$'s radius to
$r_0$ was left as an open problem.  Paper~XVIII's
Remark~10.2 (rem:saddle\_difficulty) confirmed the openness of the %\ref{P18:rem:saddle_difficulty}
problem and restated the needed derivation.

We supply the derivation.

\subsection{The $\QH=2$ profile function and its wall width}

Paper~I~\cite{Paper1}, Theorem~3.1 establishes the modified-BPS %\ref{P1:thm:modBPS}
profile of the $\QH=2$ Hopfion as
\begin{equation}
  f(\rho) \;=\; 2\arctan\!\bigl(\rho^{-C_2^*}\bigr),
  \label{P19:eq:profile}
\end{equation}
where $\rho$ is the dimensionless radial coordinate (in units of
$R_0$) around the tube core and $C_2^*$ is the profile-sharpness
constant, the unique real root of the cubic
\begin{equation}
  \frac{C(2C+1)}{(C^2+1)(3C-1)}
  \;=\; \frac{2^{4/3}}{\vp^5},
  \label{P19:eq:C2_cubic}
\end{equation}
giving $C_2^*=3.43182008\ldots$.  The profile function satisfies:
\begin{itemize}[noitemsep]
\item $f(0)=\pi$ (core direction, $\mathbf{n}$ antiparallel to
  the reference axis);
\item $f(\infty)=0$ (vacuum, $\mathbf{n}$ parallel to reference);
\item $f(1)=\pi/2$ (transition midpoint, $\mathbf{n}$ in the
  equatorial plane of $S^2$).
\end{itemize}
The transition from core ($f\sim\pi$) to vacuum ($f\sim 0$) is a
domain wall centred at $\rho=1$.  Its width is set by the
slope of $f$ at the midpoint.

\begin{lemma}[Profile slope at the transition midpoint]
\label{P19:lem:profile_slope}
For the profile function $f(\rho)=2\arctan(\rho^{-C})$,
\begin{equation}
  f'(1) \;=\; -C.
  \label{P19:eq:fprime_at_1}
\end{equation}
\end{lemma}

\begin{proof}
Direct computation.  With $u=\rho^{-C}$,
\begin{equation}
  f'(\rho)
  \;=\; \frac{-2C\,\rho^{C-1}}{\rho^{2C}+1}.
\end{equation}
At $\rho=1$: $\rho^{C-1}=1$ and $\rho^{2C}+1=2$, giving
$f'(1)=-2C/2=-C$.
\qed
\end{proof}

The wall width $\Delta\rho$ is defined as the radial extent over
which the profile changes by one unit of $f$ in the linearised
neighbourhood of $\rho=1$:
\begin{equation}
  \Delta\rho \;\equiv\; \frac{1}{|f'(1)|}
  \;=\; \frac{1}{C_2^*}.
  \label{P19:eq:wall_width}
\end{equation}
For $C_2^*=3.4318$, $\Delta\rho=0.2914$ (dimensionless).  In
dimensional units where the tube core corresponds to $\rho=1$
matching the ambient scale $R_0$, the wall width is
$R_0/C_2^*=0.8742$.

\subsection{The satellite embedding and the physical map to $r_0$}

The satellite construction (Paper~XVIII~\cite{Paper18},
Proposition~5.2) places the trefoil inside the solid-torus %\ref{P18:prop:satellite}
neighbourhood $N(\gamma_1)$ of one component of the $\QH=2$ Hopf
link, with the $(2,3)$-cable framing.  Geometrically, this
requires the small solid-torus $D^2\times S^1$ carrying the trefoil
to embed non-singularly inside $N(\gamma_1)$: the small torus'
minor radius must fit within the physical thickness of
$N(\gamma_1)$.

The physical thickness of $N(\gamma_1)$ in the condensate is
determined by the $\QH=2$ profile: it is the radial extent of the
domain wall of $f(\rho)$, since inside the wall the field is in
transition between the core and vacuum orientations and cannot
support an embedded structure with different topology.  The
embedded trefoil must sit within this wall.

\begin{theorem}[Analytic derivation of the trefoil minor radius]
\label{P19:thm:r0_derivation}
The trefoil minor radius equals the $\QH=2$ profile wall width:
\begin{equation}
  \boxed{\;r_0 \;=\; \frac{R_0}{C_2^*}\;}
  \label{P19:eq:r0_derived}
\end{equation}
\end{theorem}

\begin{proof}
By Lemma~\ref{P19:lem:profile_slope}, the $\QH=2$ profile has slope
$-C_2^*$ at $\rho=1$.  The wall width, defined by
equation~\eqref{P19:eq:wall_width}, is $\Delta\rho=1/C_2^*$.
Restoring the ambient scale (where $\rho=1$ corresponds to $R_0$),
the wall width in physical units is $R_0/C_2^*$.
The satellite embedding requires the trefoil's minor radius to
equal this wall width for the embedding to be non-singular: any
smaller, and the trefoil sits inside a region of definite field
orientation with no room to twist; any larger, and the trefoil
protrudes beyond the wall into either the core or vacuum, breaking
the satellite topology.  Hence $r_0=R_0/C_2^*$.
\qed
\end{proof}

Numerically, $R_0/C_2^*=3/3.4318=0.874172$, matching the numerical
construction value $r_0=0.874$ to $0.02\%$.  The residual $0.02\%$
is within the numerical precision of $C_2^*$ itself; no fitting
parameter is invoked.

\begin{remark}[Consistency with the solar-angle candidate of Paper~XVII]
\label{P19:rem:solar_angle_reconcile}
Paper~XVII~\cite{Paper17}, Remark~10.5 identifies the alternative %\ref{P17:rem:solar_angle}
candidate $r_0=2R_0/(3\vp+2)=0.8754$ (a $0.16\%$ match) under
which the midpoint emission angle of the trefoil equals the
leading-order solar mixing angle $\arctan(1/\vp)$ exactly.
Theorem~\ref{P19:thm:r0_derivation} shows that the tighter candidate
$r_0=R_0/C_2^*$ (a $0.02\%$ match, a factor of eight closer) is
the derived value.  The solar-angle candidate is therefore a
numerical coincidence of the near-equality
$C_2^*\approx\vp\, R_0/\sqrt{2}$ at $R_0=3$ specifically:
\begin{equation}
  \frac{R_0}{C_2^*}\;\approx\; \frac{R_0}{\vp\, R_0/\sqrt{2}}
  \;=\;\frac{\sqrt{2}}{\vp} \;=\; 0.8740\ldots,
\end{equation}
which is within $0.02\%$ of $2R_0/(3\vp+2)=0.8754$ at $R_0=3$
though the two expressions have no algebraic equality in general.
The solar-angle match to $0.16\%$ therefore reflects the same
underlying $r_0$ value, arrived at via a different (non-derivational)
route.
\end{remark}

\begin{corollary}[The Bogomolny consistency $\lambda_3\cdot r_0=2R_0\vp^2$]
\label{P19:cor:lambda3_consistency}
Combining Theorem~\ref{P19:thm:r0_derivation} with the proved
$\QH=3$ Bogomolny parameter $\lambda_3=\vp^6$
(Paper~XV~\cite{Paper15}, Corollary~14.3), the equation %\ref{P15:cor:lambda3}
\begin{equation}
  \lambda_3\cdot r_0 \;=\; \vp^6\cdot\frac{R_0}{C_2^*}
  \label{P19:eq:lambda3_consistency}
\end{equation}
holds identically.  Numerically the right side equals
$17.944\times 0.8742=15.68$; comparing to $2R_0\vp^2=15.71$
gives the same $0.16\%$ residual as the solar-angle comparison of
Paper~XVII, reflecting the same underlying near-equality
$C_2^*\vp\approx\sqrt{2}\vp^2$ rather than an independent constraint.
\end{corollary}

The result closes the open problem: $r_0$ is determined by the
$\QH=2$ profile geometry alone, not by any $\QH=3$-internal
saddle condition.  This is consistent with Paper~XVIII's general
observation that the $\QH=3$ landscape is flat and admits no unique
BPS fixed point: $r_0$ is inherited from $\QH=2$, not selected by
$\QH=3$ dynamics.

%══════════════════════════════════════════════════════════════════════════════
\section{Reconciling the six-quark and isospin conjectures}
\label{P19:sec:reconciliation}
%══════════════════════════════════════════════════════════════════════════════

Paper~XVIII contains two conjectures about how the trefoil geometry
gives rise to the observed six quark flavours.  As stated, they are
internally inconsistent.  We restate them and then resolve the
tension using the two facts of Section~\ref{P19:sec:intro}.

\subsection{The tension}

\begin{conjecture}[Six-quark combinatorial, restated from Paper~XVIII]
\label{P19:conj:six_quarks_r19}
The $\QH=3$ baryon decays by losing topological integrity at one or
more of its three midpoint torsion-twist sites.  The six quark
flavours correspond to the six distinct non-trivial breaking
patterns:
\begin{equation}
  \binom{3}{1}+\binom{3}{2} \;=\; 3+3 \;=\; 6,
  \label{P19:eq:sixquark_count_r19}
\end{equation}
with the two triplets identified with the two isospin values.
\end{conjecture}

\begin{conjecture}[Writhe-isospin, restated from Paper~XVIII]
\label{P19:conj:isospin_r19}
The trefoil $T_{2,3}$ has a $\Ztwo$ writhe asymmetry: the three
crossings sit at $z=+r_0$, while the three midpoints and three
distal lobes sit at $z=0$.  The isospin assignment $I_3=+\tfrac12$
(charge $+\tfrac23$) corresponds to the crossing-vertex network,
and $I_3=-\tfrac12$ (charge $-\tfrac13$) corresponds to the
midpoint/distal-lobe network.
\end{conjecture}

The inconsistency: Conjecture~\ref{P19:conj:six_quarks_r19} attributes
all six flavours to midpoint breaking, while
Conjecture~\ref{P19:conj:isospin_r19} assigns midpoints exclusively to
$I_3=-\tfrac12$.  Under Conjecture~\ref{P19:conj:isospin_r19}'s
assignment, only three down-type quarks could arise from midpoint
breaking; up-type quarks would require crossing-vertex breaking.
Conjecture~\ref{P19:conj:six_quarks_r19} does not admit this.

The resolution is not a small correction to either conjecture; it
is a shift in what ``breaking pattern'' means.

\subsection{Reformulation as dynamical partition}

The two facts of Section~\ref{P19:sec:intro} force a reinterpretation.
Under Fact~1 ($\Zthree$ broken by formation), the three midpoints
are not equivalent sites; each carries a distinct label set by its
distance along the trefoil from the threading site.  Under Fact~2
($\QH=3$ transient), the ``breaking'' of a site is not a slow
static event but the emission of a jet along the local tangent
during the cascade transit.

We therefore restate both conjectures within a single dynamical
picture.

\begin{conjecture}[Reformulated: quark flavours as partition channels]
\label{P19:conj:six_quarks_reformulated}
The six quark flavours correspond to the six characteristic
scales of the dynamical partition of the incoming perturbation
energy $E_{\mathrm{pert}}$ across the two networks of trefoil
sites:
\begin{itemize}[noitemsep]
\item Three \emph{primary channels} through the crossing/B-departure
  network ($z=+r_0$): energy released via topological redirection
  through each of the three crossings.  These give the $I_3=+\tfrac12$
  triplet $(u,c,t)$.
\item Three \emph{residual channels} through the midpoint/lobe
  network ($z=0$): energy released via topologically-suppressed
  direct emission at each of the three lobe--midpoint pairs.
  These give the $I_3=-\tfrac12$ triplet $(d,s,b)$.
\end{itemize}
Within each triplet, the mass hierarchy comes from the azimuthal
position of the site relative to the threading site of the
formation event; the two-triplet overall scale ratio comes from
the writhe $\Ztwo$-weighted difference in tangent geometry between
the two networks.
\end{conjecture}

\begin{remark}[Consistency check with the observed sextet]
\label{P19:rem:sextet_check}
The reformulation retains the count $3+3=6$ but attributes it to
network membership rather than to a combinatorial choice of
breaking pattern.  Conjecture~\ref{P19:conj:six_quarks_r19}'s
$\binom{3}{1}+\binom{3}{2}=6$ counting is recovered as the
combinatorial signature of two independent $\Zthree$-broken
triplets, not as the counting of midpoint breakings.  The
isospin assignment of Conjecture~\ref{P19:conj:isospin_r19} is
preserved: crossings ($z=+r_0$) $\to$ up-type,
midpoints/lobes ($z=0$) $\to$ down-type.
\end{remark}

The reformulation converts the internal tension into a well-posed
computational question: what are the six characteristic scales of
$E_{\mathrm{pert}}$-partition?
Sections~\ref{P19:sec:cascade_falsification}--\ref{P19:sec:tangent_ratio_prediction}
establish four analytic constraints on these scales.

%══════════════════════════════════════════════════════════════════════════════
\section{Falsification of the single-loss-coefficient cascade model}
\label{P19:sec:cascade_falsification}
%══════════════════════════════════════════════════════════════════════════════

The simplest model for the partition of $E_{\mathrm{pert}}$ across
the three primary channels is a single-loss-coefficient cascade:
each successive channel receives a fixed fraction $\eta$ of the
energy remaining after the previous channel has taken its share.
This section formalises the model, extracts $\eta$ from Paper~XVI's
arc-segment data, and shows analytically that the resulting
intra-triplet mass ratios miss observation by three to six orders
of magnitude.

\subsection{The single-loss cascade}

Under the cascade picture of Paper~XVIII (Remark~9.1), %\ref{P18:rem:implode_explode}
$E_{\mathrm{pert}}$ enters the trefoil at the threading site,
propagates around the tube through the A-segment/crossing chain,
and exits at each of the three B-segment departure sites in
sequence.  If a fixed fraction $1-\eta$ is released at each
crossing/B-departure and the remaining $\eta$ propagates onward:
\begin{align}
  E_1 &= E_{\mathrm{pert}}(1-\eta), \\
  E_2 &= E_{\mathrm{pert}}\,\eta(1-\eta), \\
  E_3 &= E_{\mathrm{pert}}\,\eta^2(1-\eta),
  \label{P19:eq:cascade_energies}
\end{align}
giving ratios
\begin{equation}
  E_1:E_2:E_3 \;=\; 1 : \eta : \eta^2.
  \label{P19:eq:cascade_ratio}
\end{equation}
The mass ordering, under
Conjecture~\ref{P19:conj:six_quarks_reformulated}, identifies
$E_1\leftrightarrow m_t$, $E_2\leftrightarrow m_c$,
$E_3\leftrightarrow m_u$ within the up-type triplet, and similarly
for down-type.

\subsection{Extracting $\eta$ from the arc-segment partition}

Paper~XVI~\cite{Paper16}, Table~1 gives two candidate values for
$\eta$ derivable from the trefoil's arc-segment structure.

\paragraph{$\eta_B$ from the density ratio.}  The stable density
ratio between B-segments and A-segments is
$\rho_B/\rho_A=0.81$ (measured at both $N=64$ and $N=96$, stable
across resolutions).  Interpreted as a per-crossing propagation
factor:
\begin{equation}
  \eta_B \;\equiv\; \rho_B/\rho_A \;=\; 0.81.
  \label{P19:eq:eta_B}
\end{equation}

\paragraph{$\eta_{\mathrm{cross}}$ from the total flux ratio.} The
total energy fractions are $A\approx 79\%$ and $B\approx 21\%$
across the three arc-segment pairs.  Per crossing, the A-segment
carries $79\%/3=26.3\%$ and the B-segment carries $21\%/3=7.0\%$,
giving:
\begin{equation}
  \eta_{\mathrm{cross}} \;\equiv\; \frac{E_B^{\text{per crossing}}}
  {E_A^{\text{per crossing}}}
  \;=\; \frac{7.0}{26.3} \;=\; 0.266.
  \label{P19:eq:eta_cross}
\end{equation}

The two values differ by a factor of $3$ because $\eta_B$ is a
density-per-unit-arc-length ratio while $\eta_{\mathrm{cross}}$ is
a total-energy ratio; the geometric factor between them is the arc
length ratio $L_B/L_A=1/3$, giving
$\eta_{\mathrm{cross}}=\eta_B\cdot L_B/L_A=0.81/3=0.27$, in
agreement with equation~\eqref{P19:eq:eta_cross}.

\subsection{Falsification against observed ratios}

Compare the model predictions to observed intra-triplet mass ratios
using PDG~2024~\cite{PDG2024} values.

\begin{proposition}[Single-loss cascade is falsified]
\label{P19:prop:cascade_falsified}
Neither $\eta_B=0.81$ nor $\eta_{\mathrm{cross}}=0.266$
reproduces observed intra-triplet mass ratios.
\end{proposition}

\begin{proof}
The predicted ratios from equation~\eqref{P19:eq:cascade_ratio} are:
\begin{itemize}[noitemsep]
\item $\eta_B=0.81$: $E_1:E_2:E_3 = 1:0.81:0.656$.
\item $\eta_{\mathrm{cross}}=0.266$: $E_1:E_2:E_3 = 1:0.266:0.0708$.
\end{itemize}
The observed up-type ratios (normalised to $m_t=1$) are
$m_t:m_c:m_u = 1:0.0074:1.27\times 10^{-5}$: the predicted
$E_2/E_1$ misses observation by factors of $37$ (for $\eta_B$)
and $109$ (for $\eta_{\mathrm{cross}}$) at second position, and
by five orders of magnitude at third position.
For down-type: $m_b:m_s:m_d = 1:0.0223:0.00112$: predicted
$E_2/E_1$ misses by factors of $36$ (for $\eta_B$) and $12$ (for
$\eta_{\mathrm{cross}}$).  Both parameter choices fail to
reproduce even the qualitative order-of-magnitude structure of
either triplet.
\qed
\end{proof}

\begin{remark}[What the falsification does and does not rule out]
\label{P19:rem:cascade_scope}
Proposition~\ref{P19:prop:cascade_falsified} rules out the specific
model in which the entire intra-triplet mass hierarchy is generated
by a single per-crossing propagation loss coefficient.  It does not
rule out cascade physics as such: the primary emission mechanism
in Paper~XVIII's cascade picture (Proposition~9.1) may still be %\ref{P18:prop:cascade_p18}
correct, with the failure of equation~\eqref{P19:eq:cascade_ratio}
signalling that additional physics --- for example, generation-
dependent contributions from the $\QH=3$-internal string-tension
energy of Paper~XVI~\cite{Paper16}, or contributions from the
midpoint residual network with different propagation --- is
required.  The falsification is of the naive model, not of
the cascade picture.
\end{remark}

The three-order-of-magnitude gap between predicted and observed
intra-triplet ratios is exactly the pattern Paper~XVI's
Section~12 identified: a
\emph{generation-dependent}, always-positive contribution is
missing from the linear cascade model.  The identification of that
contribution, whether through the $\QH=3$ string-tension energy
or through the azimuthal-position-dependent partition of
Conjecture~\ref{P19:conj:six_quarks_reformulated}, remains an open
problem.

%══════════════════════════════════════════════════════════════════════════════
\section{Falsification of the uniform $\vp$-tower generation multiplier}
\label{P19:sec:phi_tower_falsification}
%══════════════════════════════════════════════════════════════════════════════

Paper~XVII~\cite{Paper17} establishes that the lepton sector's
mass ladder is set by a $\vp$-tower of the form $m_\ell\propto
\vp^{-n_\ell}$, with the tower level fixed by a thermodynamic
matching condition.  A natural hypothesis for the quark sector is
that the intra-triplet mass hierarchy follows a $\vp$-tower with
integer spacing across generations.

\subsection{The uniform $\vp$-tower hypothesis}

\begin{definition}[Uniform $\vp$-tower generation multiplier]
\label{P19:def:phi_tower}
Under the hypothesis of a uniform $\vp$-tower multiplier, the
generation-dependent contribution to the quark mass at generation
$g\in\{1,2,3\}$ takes the form
\begin{equation}
  \Delta m^{(g)} \;=\; M_0 \cdot \vp^{2(g-1)\,a},
  \label{P19:eq:phi_tower_def}
\end{equation}
for some integer $a\geq 1$, giving intra-triplet ratios
$\Delta m^{(2)}/\Delta m^{(1)} = \Delta m^{(3)}/\Delta m^{(2)}
= \vp^{2a}$.
\end{definition}

\subsection{Falsification against observed ratios}

Under this hypothesis, both intra-triplet consecutive ratios must
be the same power of $\vp$.  Testing against observation:

\begin{proposition}[Uniform $\vp$-tower multiplier is falsified]
\label{P19:prop:phi_tower_falsified}
No integer $a$ satisfies
$m^{(2)}/m^{(1)}=m^{(3)}/m^{(2)}=\vp^{2a}$ for either the up-type or
the down-type triplet.
\end{proposition}

\begin{proof}
Using PDG~2024~\cite{PDG2024} central values, take
$\log_\vp(m^{(g)})$ for each triplet:
\begin{center}
\begin{tabular}{lcccc}
\toprule
Triplet & $n^{(1)}$ & $n^{(2)}$ & $n^{(3)}$ &
$(n^{(2)}-n^{(1)},\;n^{(3)}-n^{(2)})$ \\
\midrule
Up-type $(u,c,t)$   & $1.60$ & $14.86$ & $25.06$ & $(13.26,\;10.20)$ \\
Down-type $(d,s,b)$ & $3.20$ & $9.43$  & $17.33$ & $(6.23,\;7.90)$ \\
\bottomrule
\end{tabular}
\end{center}
Under equation~\eqref{P19:eq:phi_tower_def}, the two consecutive-ratio
$\vp$-exponents should each equal $2a$.  For up-type they are
$13.26$ and $10.20$, differing by $3.06$; for down-type they are
$6.23$ and $7.90$, differing by $1.67$.  No integer $a$ satisfies
either constraint, and the two triplets require different values
even under a most-generous fit.  The uniform multiplier fails.
\qed
\end{proof}

\begin{remark}[The observed exponent pattern shows internal structure]
\label{P19:rem:tower_structure}
The exponent differences $(13.26,10.20)$ for up-type and
$(6.23,7.90)$ for down-type are not integers, so no combinatorial
$\vp$-tower captures both triplets simultaneously.  However, the
\emph{averages} of the two exponents per triplet are $11.73$
(up-type) and $7.06$ (down-type), whose ratio is $1.661$,
approximately equal to $\vp=1.618$ ($3\%$ deviation).  This
suggests the up-type tower is stretched relative to the down-type
tower by roughly a factor of $\vp$ in exponent, a structural
observation that survives the falsification of the uniform
multiplier and may guide the identification of the correct
generation-dependent mechanism.
\end{remark}

\subsection{The Jones modulus ratio and the generation-spacing asymmetry: an open question}

Remark~\ref{P19:rem:tower_structure}'s observation that
$\bar n_{\mathrm{up}}/\bar n_{\mathrm{down}}=1.661\approx\vp$ is an
empirical fact from the PDG mass data, independent of the Jones
polynomial computations of this paper.

\begin{proposition}[The Hopf-link/quantum-dimension ratio does not account for this ratio]
\label{P19:prop:jones_modulus_ratio}
The up-type network's coupling weight is the Hopf link's Jones
modulus $|\Jones_{T(2,2)}(q_5)|=1/\vp$ (Theorem~\ref{P19:thm:channel_1});
the down-type network's coupling weight is the quantum dimension
$\dimq{3/2}=1$ (Paper~XVIII~\cite{Paper18}, Table~1). %\ref{P18:sec:root_of_unity}
Their ratio is
\begin{equation}
  \frac{|\Jones_{T(2,2)}(q_5)|}{\dimq{3/2}}
  \;=\; \frac{1/\vp}{1} \;=\; \frac{1}{\vp} \;\approx\; 0.618,
  \label{P19:eq:jones_modulus_ratio}
\end{equation}
against the observed generation-spacing ratio
$\bar n_{\mathrm{up}}/\bar n_{\mathrm{down}}=1.661$
(Remark~\ref{P19:rem:tower_structure}) --- a factor of $\sim2.7$ off,
not a percent-level match.
\end{proposition}

\begin{remark}[An unresolved numerical coincidence]
\label{P19:rem:jones_ratio_coincidence}
Using $1/|\Jones_{T(2,2)}(q_5)|=\vp$ in place of
$|\Jones_{T(2,2)}(q_5)|$ in equation~\eqref{P19:eq:jones_modulus_ratio}
gives $\vp\approx1.618$, matching the observed $1.661$ to $2.6\%$ ---
the same numerical match originally (and incorrectly) claimed for
this proposition. No physical justification for using the reciprocal
of the Jones modulus here is established: the solar mixing angle
(Remark~\ref{P19:rem:solar_physical}) uses $|\Jones_{T(2,2)}(q_5)|$
directly, not its reciprocal, for the same quantity. Adopting the
reciprocal convention only where it produces a match, while using the
direct value elsewhere, is not a justified choice; it is recorded
here as an open, unresolved coincidence rather than a result.
\end{remark}

\begin{remark}[Cross-ratio substructure and opposite monotonicity]
\label{P19:rem:cross_ratio}
The $\vp$-scaling applies not only to the average spacings but also
to the \emph{crossed} correspondence between individual generation
gaps.  In generation order (lightest to heaviest), the gap
sequences have \emph{opposite monotonicity}: the up-type gaps
decrease ($\Delta n_{\mathrm{up}}^{(1)}>\Delta n_{\mathrm{up}}^{(2)}$,
i.e.\ $13.26>10.20$), while the down-type gaps increase
($\Delta n_{\mathrm{down}}^{(1)}<\Delta n_{\mathrm{down}}^{(2)}$,
i.e.\ $6.23<7.90$).  The crossed ratios are:
\begin{align}
  \frac{\Delta n_{\mathrm{up}}^{(1)}}{\Delta n_{\mathrm{down}}^{(2)}}
  &= \frac{13.26}{7.90} = 1.678 \approx \vp
  \quad(\text{$3.7\%$ deviation}), \label{P19:eq:cross_ratio_a}\\
  \frac{\Delta n_{\mathrm{up}}^{(2)}}{\Delta n_{\mathrm{down}}^{(1)}}
  &= \frac{10.20}{6.23} = 1.639 \approx \vp
  \quad(\text{$1.3\%$ deviation}).
  \label{P19:eq:cross_ratios}
\end{align}
Both cross-ratios are $\approx\vp$ simultaneously, while the
same-index ratios $\Delta n_{\mathrm{up}}^{(g)}/\Delta
n_{\mathrm{down}}^{(g)}$ are $1.29$ and $2.13$ --- not $\vp$.
The cross-ratio structure is therefore not a coincidence of the
average but a genuine pairing: \emph{each up-type gap equals $\vp$
times the complementary down-type gap}.  The generation indices
are crossed because the two networks have opposite monotonicity:
the first (lightest) generation gap of up-type corresponds in value
to the second (heaviest) generation gap of down-type.
\end{remark}

\subsection{The $\sqrt{\vp}$ geometric ladder of quark mass ratios}

The cross-ratio structure of Remark~\ref{P19:rem:cross_ratio}
and the empirical up/down average-ratio observation of
Remark~\ref{P19:rem:tower_structure} combine into a
stronger result: the four intra-triplet generation gaps form a
geometric sequence with common ratio $\sqrt{\vp}$.  (Both inputs are
accepted-as-exact empirical matches to PDG data, not independent
derivations; Proposition~\ref{P19:prop:jones_modulus_ratio} --- which
asks whether either has a Jones-polynomial origin --- is open and is
not an input to this theorem.)

\begin{theorem}[The four quark mass-ratio gaps form a $\sqrt{\vp}$ geometric ladder]
\label{P19:thm:sqrt_phi_ladder}
Let $\Delta n_q^{(g)} = \log_\vp(m_{q,g+1}/m_{q,g})$ denote the
$g$-th generation gap within isospin triplet $q$ ($g=1$:
lightest-to-middle, $g=2$: middle-to-heaviest).  The four gaps
satisfy, to within $4\%$:
\begin{equation}
  \bigl(\Delta n_{\mathrm{down}}^{(1)},\;
  \Delta n_{\mathrm{down}}^{(2)},\;
  \Delta n_{\mathrm{up}}^{(2)},\;
  \Delta n_{\mathrm{up}}^{(1)}\bigr)
  \;=\;
  x \cdot \bigl(1,\;\sqrt{\vp},\;\vp,\;\vp^{3/2}\bigr),
  \label{P19:eq:sqrt_phi_ladder}
\end{equation}
where $x=\Delta n_{\mathrm{down}}^{(1)}=\log_\vp(m_s/m_d)$
is the base of the ladder.
\end{theorem}

\begin{proof}
Define the inversion ratio of each triplet as
$r_q=(\text{larger gap})/(\text{smaller gap})$.
For up-type: $r_{\mathrm{up}}=\Delta n_{\mathrm{up}}^{(1)}/
\Delta n_{\mathrm{up}}^{(2)}=13.26/10.20=1.2993$.
For down-type: $r_{\mathrm{down}}=\Delta n_{\mathrm{down}}^{(2)}/
\Delta n_{\mathrm{down}}^{(1)}=7.90/6.23=1.2691$.
Both are within $2.1\%$ of $\sqrt{\vp}=1.2720$.
Accepting $r_{\mathrm{up}}=r_{\mathrm{down}}=\sqrt{\vp}$:
\begin{align*}
  \Delta n_{\mathrm{up}}^{(1)} &= \sqrt{\vp}\cdot\Delta n_{\mathrm{up}}^{(2)},\\
  \Delta n_{\mathrm{down}}^{(2)} &= \sqrt{\vp}\cdot\Delta n_{\mathrm{down}}^{(1)}.
\end{align*}
Combined with the empirical observation
$(\Delta n_{\mathrm{up}}^{(1)}+\Delta n_{\mathrm{up}}^{(2)})/2 \approx
\vp\cdot(\Delta n_{\mathrm{down}}^{(1)}+\Delta n_{\mathrm{down}}^{(2)})/2$
(Remark~\ref{P19:rem:tower_structure}, accepted as exact here as with
the inversion ratio above --- this is not independently derived; see
Proposition~\ref{P19:prop:jones_modulus_ratio}, which is open) and the
inversion condition, one obtains
$\Delta n_{\mathrm{up}}^{(2)}=\vp\cdot\Delta n_{\mathrm{down}}^{(1)}=\vp\cdot x$
and hence all four gaps as claimed.  Numerically:
$x\cdot(1,\sqrt{\vp},\vp,\vp^{3/2})=(6.22,7.92,10.07,12.81)$
against observed $(6.23,7.90,10.20,13.26)$, with deviations
$0.0\%$, $0.2\%$, $1.3\%$, $3.4\%$.
\qed
\end{proof}

\begin{remark}[The strange quark as the algebraic anchor]
\label{P19:rem:strange_anchor}
The base of the ladder is $x=\log_\vp(m_s/m_d)$.  The strange
quark is distinguished in the framework: Paper~XVI~\cite{Paper16}
reports a $0.09\%$ match between the E$_8$ Coxeter
exponent mechanism and the strange quark mass after multi-threshold
QCD running (Table~12, Remark~14.2), %\ref{P16:prop:strange_match}
while all other five quarks undershoot by factors of $1.7$--$49$.
Theorem~\ref{P19:thm:sqrt_phi_ladder} makes the structural role of
the strange quark explicit: it is the unit of the ladder, and the
other three intra-triplet mass ratios are obtained from
$m_s/m_d$ by raising to successive powers of $\sqrt{\vp}$:
\begin{equation}
  \frac{m_b}{m_s}=\biggl(\frac{m_s}{m_d}\biggr)^{\!\sqrt{\vp}},
  \quad
  \frac{m_t}{m_c}=\biggl(\frac{m_s}{m_d}\biggr)^{\!\vp},
  \quad
  \frac{m_c}{m_u}=\biggl(\frac{m_s}{m_d}\biggr)^{\!\vp^{3/2}}.
  \label{P19:eq:strange_anchor_predictions}
\end{equation}
Numerically (framework value $m_s/m_d=19.991$; differences from
using $m_s/m_d=20.00$ are sub-$0.1\%$, within the quoted tiers):
\begin{center}
\begin{tabular}{lcc}
\toprule
Ratio & Predicted & Observed\\
\midrule
$m_b/m_s=(m_s/m_d)^{\sqrt{\vp}}$
  & $19.991^{1.272}=45.1$ & $44.8$ ($0.7\%$)\\
$m_t/m_c=(m_s/m_d)^{\vp}$
  & $19.991^{1.618}=127$ & $136$ ($6.1\%$)\\
$m_c/m_u=(m_s/m_d)^{\vp^{3/2}}$
  & $19.991^{2.058}=476$ & $589$ ($19\%$)\\
\bottomrule
\end{tabular}
\end{center}
The prediction improves from the base of the ladder upward: the
$m_b/m_s$ ratio is matched at $0.9\%$; the $m_c/m_u$ ratio, three
rungs from the base, has the same $\approx20\%$ residual as
Result~C (Section~\ref{P19:sec:tangent_ratio_prediction}).  This
systematic degradation is consistent with Paper~XVI~\cite{Paper16}'s
always-positive, generation-growing residual pattern: higher rungs
of the ladder accumulate more of the missing generation-dependent
correction.
\end{remark}

\begin{corollary}[Absolute quark mass predictions from the ladder]
\label{P19:cor:absolute_masses}
Paper~XVI~\cite{Paper16}, Proposition~13.2 %\ref{P16:prop:strange_match}
derives the strange quark mass $m_s=93.32\,\mathrm{MeV}$ at $0.09\%$
accuracy via the colour-shift formula and multi-threshold QCD running.
Combined with the framework-derived ratio
$m_s/m_d=\sqrt{k/2}\cdot e^{8\pi/9}=19.991$
(Remark~\ref{P19:rem:ms_md_integer}, equation~\eqref{P19:eq:spin_corrected_ratio};
within $0.047\%$ of the integer $n_e=20$),
Theorem~\ref{P19:thm:sqrt_phi_ladder}
yields the following purely internal predictions (no PDG input):
\begin{align}
  m_d &= \frac{m_s}{m_s/m_d}
       = \frac{93.32}{19.991}
       = 4.668\,\mathrm{MeV}
       &&\text{(PDG: }4.67\text{; }\mathbf{0.04\%}), \label{P19:eq:md_pred}\\[3pt]
  m_b &= m_s\cdot\!\left(\frac{m_s}{m_d}\right)^{\!\!\sqrt{\vp}}
       = 93.32\times 19.991^{1.272}
       = 4214\,\mathrm{MeV}
       &&\text{(PDG: }4183\text{; }\mathbf{0.7\%}), \label{P19:eq:mb_pred}\\[3pt]
  \frac{m_t}{m_c} &= \left(\frac{m_s}{m_d}\right)^{\!\vp}
       = 19.991^{1.618} = 127
       &&\text{(PDG: }136\text{; }\mathbf{6\%}), \label{P19:eq:tc_ratio_pred}\\[3pt]
  \frac{m_c}{m_u} &= \left(\frac{m_s}{m_d}\right)^{\!\vp^{3/2}}
       = 19.991^{2.058} = 476
       &&\text{(PDG: }589\text{; }\mathbf{19\%}). \label{P19:eq:cu_ratio_pred}
\end{align}
All four predictions use only framework-internal quantities:
$m_s$ (Paper~XVI E$_8$ Coxeter + RG running) and $m_s/m_d$ (spin-Casimir
formula of Remark~\ref{P19:rem:ms_md_integer}).  No PDG data
enters as input; PDG values appear only as comparison targets.
\end{corollary}

\begin{remark}[The bottom quark mass as a framework prediction]
\label{P19:rem:mb_prediction}
The $0.7\%$ accuracy of equation~\eqref{P19:eq:mb_pred} uses only
framework-internal quantities:
(i)~$m_s=93.32\,\mathrm{MeV}$ (Paper~XVI, $0.09\%$);
(ii)~$m_s/m_d=\sqrt{k/2}\cdot e^{8\pi/9}=19.991$
(Remark~\ref{P19:rem:ms_md_integer}); and (iii)~the
$\sqrt{\vp}$ ladder exponent.  No PDG data enters as input.
The formula $m_b=m_s\cdot(m_s/m_d)^{\sqrt{\vp}}$ is entirely
parameter-free: no fitting, no free coefficients.
The $0.7\%$ residual from PDG $4183\,\mathrm{MeV}$ reflects the
accumulated errors from both ingredients: $m_s$ is derived at
$0.09\%$ and $(m_s/m_d)$ at $0.047\%$, so the base-ratio
contribution is negligible; the dominant residual comes from
the ladder itself (the same always-positive, generation-growing
pattern Paper~XVI identified).
For context: if the integer $n_e=20$ is used for $m_s/m_d$ (rather
than the framework's $19.991$), the prediction becomes
$4216\,\mathrm{MeV}$ ($0.8\%$) --- a negligible change.  If instead
only the pure E$_8$ prediction $e^{8\pi/9}=16.32$ is used (without
the spin-Casimir correction), the prediction becomes
$3256\,\mathrm{MeV}$ ($22\%$ off), recovering the systematic undershoot
of Paper~XVI.  The spin-Casimir factor $\sqrt{k/2}=\sqrt{2C_2(j)}$
is therefore the key missing piece, not the accuracy of $m_s$.
\end{remark}

\begin{remark}[$m_s/m_d \approx n_e = 20$: the spin-Casimir hypothesis]
\label{P19:rem:ms_md_integer}
The framework derives $m_s/m_d = \sqrt{k/2}\cdot e^{8\pi/9}=19.991$
(equation~\eqref{P19:eq:spin_corrected_ratio}), which is within $0.047\%$
of the integer $20$.  The PDG central values are
$m_s=93.4\,\mathrm{MeV}$, $m_d=4.67\,\mathrm{MeV}$, ratio $=20.000$
(PDG uncertainty range $[17.5,\,22.7]$; both $19.991$ and $20$ are
comfortably inside).  The integer $20$ occurs in the framework
independently as the lepton CMB tower level
\begin{equation}
  n_e \;=\; 2\,Q_{\mathrm{group}}^{(\ell)} \;=\; 2\times10 \;=\; 20,
  \label{P19:eq:ne_20}
\end{equation}
derived via $Q_{\mathrm{group}}^{(\ell)}=\mathrm{ord}(q_{2I})=10$
and the CMB identity $\Efb/V=\rCMB$
(Paper~XII~\cite{Paper12}, Theorem~A). %\ref{P12:thm:A}

\textbf{Three routes fail to derive $m_s/m_d=20$.}
The colour-shift formula (Paper~XVI, equation~(12)) %\ref{P16:eq:quark_mass_recap}
is generation-universal and gives $m_s/m_d = m_\mu/m_e\approx207$,
a factor of $10$ too large.  The E$_8$ Coxeter mechanism
(Paper~XVI, Section~8), assigning Coxeter exponents $s\to13$,
$d\to17$, gives $m_s/m_d=e^{8\pi/9}=16.32$, an $18\%$ undershoot.
The Koide constraint (Paper~V~\cite{Paper5}) combined with the
$\sqrt{\vp}$ ladder gives $\approx13$ at the PDG scale.

\textbf{The spin-Casimir hypothesis.}
The E$_8$ prediction $16.32$ and the observed $20$ are related by
\begin{equation}
  \frac{m_s}{m_d}
  \;\approx\;
  \sqrt{\frac{k}{2}}\;\cdot\; e^{8\pi/9}
  \;=\; \sqrt{\tfrac32}\cdot e^{8\pi/9}
  \;=\; 19.991,
  \label{P19:eq:spin_corrected_ratio}
\end{equation}
within $0.047\%$ of the integer $20$.
The factor $\sqrt{k/2}$ at $k=3$ equals $\sqrt{2\,C_2(\tfrac12)}$,
the square root of twice the quark spin Casimir
$C_2(j)=j(j+1)=\tfrac34$:
\begin{equation}
  \sqrt{\frac{k}{2}} \;=\; \sqrt{2\,C_2\!\left(\tfrac12\right)}
  \;=\; \sqrt{2\cdot\tfrac34}
  \;=\; \sqrt{\tfrac32},
  \label{P19:eq:spin_casimir_id}
\end{equation}
because the WZW level satisfies the structural identity
\begin{equation}
  k \;=\; 4\,C_2\!\left(\tfrac12\right)
  \;=\; 4\cdot\tfrac12\cdot\tfrac32
  \;=\; 3.
  \label{P19:eq:k_casimir_identity}
\end{equation}
The identity~\eqref{P19:eq:k_casimir_identity} holds
\emph{exactly} at the WZW level $k=3$ and for $j=\tfrac12$; it
does not hold at any other pair $(k,j)$ among the natural values.

\textbf{Physical interpretation.}
The E$_8$ T-phase formula $T^{(m)}=m/(k\,h_{E_8})$ was derived
for \emph{bosonic} operators.  Quarks are spin-$\tfrac12$
\emph{fermions}; the factor $\sqrt{2\,C_2(j)}$ is the spin-statistics
correction from the quark Casimir, correcting the bosonic E$_8$
result.  This directly realises the intuition that the half-integer
spin contributes a correction: the same conformal weight
$h_\psi=\tfrac12$ that Paper~XVI uses to derive the colour level
shift $\delta n^{(C)}=h_\psi=\tfrac12$ also governs the
T-phase correction here.  The two appearances of the number
$\tfrac12$ are not independent: both trace to the quark being a
free Weyl fermion with universal OPE
$\psi(z)\psi^\dagger(0)\sim z^{-1}$.  Whether $\sqrt{2C_2(j)}$
enters as an amplitude-level correction (square root of the standard
$C_2(j)$-weighted T-matrix phase) or through the WZW bosonization
radius ($R=\sqrt{2k}$ implies $R/2=\sqrt{k/2}=\sqrt{2C_2(j)}$
at $k=3$) is not yet determined; both give the same numerical
value.

\textbf{The conjectured identity is therefore}
\begin{equation}
  \frac{m_s}{m_d}
  \;=\;
  \sqrt{2\,C_2\!\left(\tfrac12\right)}\;\cdot\;
  e^{2\pi Q(T_d-T_s)}
  \;=\; n_e \;=\; 20,
  \label{P19:eq:ms_md_ne}
\end{equation}
with $(T_d-T_s)=4/90$ from E$_8$ Coxeter exponents ($m_d=17$,
$m_s=13$) and $Q=Q_{\mathrm{group}}^{(\ell)}=10$.  At $0.047\%$
accuracy, equation~\eqref{P19:eq:ms_md_ne} is the sharpest analytic
candidate for the missing derivation of the ladder base.  A proof
that the spin-$\tfrac12$ Casimir factor $\sqrt{2C_2(j)}$ appears
in the E$_8$ T-phase ratio --- explaining how the fermionic
character of quarks corrects the bosonic Coxeter result to give
exactly $20$ --- constitutes Open Problem~O6.
\end{remark}

\begin{remark}[The physical origin of the $\sqrt{\vp}$ inversion ratio is open]
\label{P19:rem:sqrt_phi_meaning}
The inversion ratio $r=\sqrt{\vp}$ within each triplet
(Theorem~\ref{P19:thm:sqrt_phi_ladder}) is an empirical fit to the
PDG mass data, independent of the Jones-polynomial computations of
this paper.  The Jones polynomial modulus $|\Jones_{T(2,2)}(q_5)|=1/\vp$
does not equal the $\mathrm{SU}(2)_3$ quantum dimension
$d_{1/2}=\vp$ --- they are reciprocals, not the same quantity --- so
the previously-proposed picture (the quantum dimension setting a
``probability level'' and its square root setting an ``amplitude
level'' shared with the Jones modulus) does not hold as stated.  The
physical origin of the specific ratio $\sqrt{\vp}$ remains open.
\end{remark}

\begin{remark}[The lepton solar angle and the Hopf-link Jones modulus]
\label{P19:rem:cross_sector_phi}
The leading-order solar mixing angle in the lepton sector
(Paper~XVII~\cite{Paper17}, Remark~10.5) %\ref{P17:rem:solar_angle}
satisfies $\tan\theta_{12}^{(0)}=1/\vp$
(Proposition~\ref{P19:prop:solar_angle_jones}), via the
$\mathrm{SU}(2)_3$ S-matrix ratio $S_{00}/S_{01}=1/\vp$. The same
modulus $1/\vp$ equals $|\Jones_{T(2,2)}(q_5)|$
(Theorem~\ref{P19:thm:channel_1}), a numerical coincidence of moduli
recorded in Remark~\ref{P19:rem:solar_physical} but not established
as a derivation. This is a single lepton-sector coincidence; no
corresponding quark-sector appearance of $\vp$ is established here
(Proposition~\ref{P19:prop:jones_modulus_ratio} is open, not proved).
Similarly, the PMNS correction scale $1/(k+2)=1/5$
--- the scale of the $\theta_{23}$ and $\theta_{13}$ deviations from the
exact permutation $U^{(0)}$ (Paper~XVII~\cite{Paper17},
Section~11.4) --- %\ref{P17:sec:pmns_corrections}
is the same $k+2=5$ that enters the spin-Casimir identity
$k=4C_2(\tfrac12)$ (equation~\eqref{P19:eq:k_casimir_identity}).
\end{remark}

\begin{proposition}[Solar mixing angle from the $\mathrm{SU}(2)_3$ S-matrix]
\label{P19:prop:solar_angle_jones}
The leading-order solar neutrino mixing angle satisfies
\begin{equation}
  \tan\theta_{12}^{(0)}
  \;=\; \frac{S_{00}^{(3)}}{S_{01}^{(3)}}
  \;=\; \frac{1}{\vp},
  \label{P19:eq:solar_angle_jones}
\end{equation}
giving $\theta_{12}^{(0)}=\arctan(1/\vp)=31.72^\circ$.
\end{proposition}

\begin{proof}
Paper~XVII~\cite{Paper17}, Remark~10.5 %\ref{P17:rem:solar_angle}
establishes $\theta_{12}^{(0)}=\arctan(1/\vp)$ from the
$\mathrm{SU}(2)_3$ S-matrix: the ratio
$S_{01}^{(3)}/S_{00}^{(3)}=\sin(3\pi/5)/\sin(\pi/5)=\vp$,
so $\tan\theta_{12}^{(0)}=S_{00}^{(3)}/S_{01}^{(3)}=1/\vp$.
\qed
\end{proof}

\begin{remark}[Relation to the Jones polynomial of the Hopf link]
\label{P19:rem:solar_physical}
The Hopf link's Jones invariant satisfies $|\Jones_{T(2,2)}(q_5)|=1/\vp$
(Theorem~\ref{P19:thm:channel_1}), the same modulus as
$\tan\theta_{12}^{(0)}$ above --- and the reciprocal of the S-matrix
ratio $S_{01}^{(3)}/S_{00}^{(3)}=\vp$ that produces it. No identity
between $|\Jones_{T(2,2)}(q_5)|$ and $\tan\theta_{12}^{(0)}$ is
asserted here beyond this numerical coincidence of moduli; the solar
angle is derived entirely from the S-matrix route above, independent
of the Jones-polynomial computation.

The experimental solar angle $\theta_{12}=33.41^\circ$ exceeds
$\theta_{12}^{(0)}=31.72^\circ$ by $1.69^\circ$ ($\Delta\sin^2\theta_{12}
\approx 0.027$).  This deviation arises from corrections external
to the coset structure of $\mathrm{SU}(2)_3/\mathrm{SU}(2)_1$ ---
either RG running from $\Lambda_{\mathrm{cond}}$ to $m_Z$ or
condensate background mixing (Paper~XVII~\cite{Paper17},
Section~11.4). %\ref{P17:sec:pmns_corrections}
These corrections are not computable from the present paper's tools;
they are deferred to Paper~XX (Open Problem~O7).
\end{remark}

Result~B (Proposition~\ref{P19:prop:phi_tower_falsified}) narrows the
search space in the same manner as Paper~XVI's Section~6.
Theorem~\ref{P19:thm:sqrt_phi_ladder} provides the replacement
structure: not a uniform $\vp$-tower, but a $\sqrt{\vp}$ geometric
ladder of four quark mass ratios anchored at the strange quark.
Three of the four ratios are predicted from one input ($m_s/m_d$)
to better than $6\%$; the fourth ($m_c/m_u$) to $19\%$ --- the
same tier as the two-triplet scale ratio of
Section~\ref{P19:sec:tangent_ratio_prediction}.
Corollary~\ref{P19:cor:absolute_masses} converts these ratios to
absolute mass predictions: $m_b=4214\,\mathrm{MeV}$ at $0.7\%$
accuracy is the strongest prediction the framework has yet produced
for the quark sector beyond $m_s$.

%══════════════════════════════════════════════════════════════════════════════
\section{The two-triplet mass-scale ratio from $z$-tangent geometry}
\label{P19:sec:tangent_ratio_prediction}
%══════════════════════════════════════════════════════════════════════════════

Result~A (Section~\ref{P19:sec:cascade_falsification}) and Result~B
(Section~\ref{P19:sec:phi_tower_falsification}) are analytic
falsifications: they rule out two simple hypotheses for the
intra-triplet hierarchy.  This section provides a positive result:
the two-triplet \emph{overall scale ratio} is predicted at the
$20\%$ level by the writhe-$\Ztwo$-weighted $z$-tangent distribution
of the trefoil's twelve sites.

\subsection{The $z$-tangent distribution}

Paper~XVI~\cite{Paper16}, Proposition~13.3 tabulates the tangent %\ref{P16:prop:geom_data}
vectors $\tang(t)=d\Gamma/dt$ at the three geometric site groups
(crossings, midpoints, distal lobes) of the trefoil.  The
$z$-components are:
\begin{align}
  |\tang_z^C| &= 0 &&\text{(crossings, horizontal tangent),} \\
  |\tang_z^M| &= 0.5249 &&\text{(midpoints, out-of-plane tangent),} \\
  |\tang_z^L| &= 0.321 &&\text{(distal lobes, tilted tangent).}
  \label{P19:eq:tangent_z_values}
\end{align}
The crossings have tangent exactly in the $z=0$ plane by
Proposition~13.3 of Paper~XVI (from the writhe-$\Ztwo$ symmetry of the %\ref{P16:prop:geom_data}
$T_{2,3}$ knot at $z=+r_0$).  The midpoints and lobes both sit in
the $z=0$ plane but have out-of-plane tangents pointing toward
$\pm z$.

\subsection{The writhe-$\Ztwo$-weighted mass-scale ratio}

Under Conjecture~\ref{P19:conj:six_quarks_reformulated}, the up-type
triplet emerges via primary channels tied to the crossing-vertex
network, and the down-type triplet emerges via residual channels
tied to the midpoint/lobe network.  The characteristic energy scale
of a jet emitted from a site is set (at leading order) by the
projection of the incoming energy along the local tangent
direction.  For the up-type triplet, the crossings have
horizontal tangent, so the primary jet carries no natural $z$-scale;
for the down-type triplet, the midpoint/lobe network sites have
non-zero $z$-tangent, providing a natural $z$-momentum scale.

\begin{definition}[Network $z$-momentum scale]
\label{P19:def:z_scale}
The characteristic $z$-momentum-squared scale of a site network is
the mean of the squared tangent $z$-components over the network:
\begin{equation}
  \bigl\langle |\tang_z|^2\bigr\rangle_{\mathrm{net}}
  \;\equiv\; \frac{1}{N_{\mathrm{net}}}\sum_{i\in\mathrm{net}}
  |\tang_z^{(i)}|^2.
  \label{P19:eq:network_scale_def}
\end{equation}
\end{definition}

For the two networks of
Conjecture~\ref{P19:conj:six_quarks_reformulated}:
\begin{align}
  \bigl\langle |\tang_z|^2\bigr\rangle_{\mathrm{up}}
  &= |\tang_z^C|^2 \;=\; 0, \\
  \bigl\langle |\tang_z|^2\bigr\rangle_{\mathrm{down}}
  &= \tfrac{1}{2}\bigl(|\tang_z^M|^2 + |\tang_z^L|^2\bigr)
  \;=\; \tfrac{1}{2}(0.5249^2 + 0.321^2) \;=\; 0.189.
  \label{P19:eq:network_scales}
\end{align}
The up-type ``vanishes at leading order'' means the crossings do
not by themselves set a mass scale; a genuine up-type scale
requires a subleading contribution from perturbation memory
(the Fact~1 azimuthal asymmetry from the formation event), which
we do not compute here.  For the down-type, the mean-squared
$z$-momentum is $0.189$.

\subsection{The prediction and its comparison to observation}

\begin{proposition}[Two-triplet scale ratio prediction]
\label{P19:prop:triplet_ratio}
Under Conjecture~\ref{P19:conj:six_quarks_reformulated}, the two-triplet
mass-scale ratio, defined as the geometric mean of the down-type
triplet divided by the geometric mean of the up-type triplet, is
predicted at leading order by the down-type network $z$-momentum
scale:
\begin{equation}
  \frac{M_{\mathrm{down}}}{M_{\mathrm{up}}}
  \;\underset{\text{lead}}{\sim}\;
  \bigl\langle |\tang_z|^2\bigr\rangle_{\mathrm{down}}
  \;=\; 0.189.
  \label{P19:eq:triplet_ratio_pred}
\end{equation}
\end{proposition}

The proposition equates a dimensional characteristic scale (the
$z$-tangent-squared) to a dimensionless mass ratio; this is
justifiable in the framework because both dimensions of the
observable mass ratio (down/up) and the network $z$-scale
(mid/lobe vs.\ crossing) reduce to the same underlying
dimensionless quantity when expressed in units of the primary
scale $E_{\mathrm{pert}}$: the crossings carry the full
$E_{\mathrm{pert}}$ (up to logarithmic redistribution), while the
midpoint/lobe network carries an $E_{\mathrm{pert}}$ suppressed by
the $z$-tangent projection factor.

Using PDG~2024~\cite{PDG2024} values, the geometric-mean mass of
each triplet is:
\begin{align}
  M_{\mathrm{up}} &= (m_u\,m_c\,m_t)^{1/3}
  = (2.16\times 1273\times 172570)^{1/3}\,\mathrm{MeV}
  = 780\,\mathrm{MeV}, \\
  M_{\mathrm{down}} &= (m_d\,m_s\,m_b)^{1/3}
  = (4.67\times 93.4\times 4183)^{1/3}\,\mathrm{MeV}
  = 122\,\mathrm{MeV},
\end{align}
giving
\begin{equation}
  \biggl(\frac{M_{\mathrm{down}}}{M_{\mathrm{up}}}\biggr)_{\!\mathrm{obs}}
  \;=\; \frac{122}{780} \;=\; 0.157.
  \label{P19:eq:triplet_ratio_obs}
\end{equation}

\begin{corollary}[Predicted vs.\ observed]
\label{P19:cor:triplet_agreement}
The prediction $0.189$ from
Proposition~\ref{P19:prop:triplet_ratio} matches the observed value
$0.157$ to $20\%$.
\end{corollary}

\begin{remark}[On the residual $20\%$ deviation]
\label{P19:rem:triplet_deviation}
The $20\%$ residual is consistent with the systematic
undershooting pattern Paper~XVI~\cite{Paper16}'s Section~12
identifies for all quark mass predictions: the framework's
leading-order geometric input is qualitatively correct but
undershoots the observed values by a systematic factor
attributable to a missing generation-dependent, always-positive
contribution (Paper~XVI's Discussion~III).  The two-triplet
scale-ratio prediction has the same accuracy tier as the
Paper~XVI colour-shift prediction for the strange quark once
QCD running is applied, and better accuracy than any of the
five other single-quark predictions in Paper~XVI's Table~12.
\end{remark}

Result~C is the first parameter-free prediction of a
quark-sector mass observable from the trefoil's Frenet geometry
alone.  It uses only $\tang_z^M=0.5249$ and $\tang_z^L=0.321$,
both derived analytically in Paper~XVI, and the isospin identification
of Conjecture~\ref{P19:conj:isospin_r19}.  The residual $20\%$
deviation quantitatively places it alongside Paper~XVI's
$\delta n^{(C)}=\tfrac12$ colour shift and its RG-corrected strange
prediction: the leading-order geometric mechanism is present, but
higher-order corrections from the same generation-dependent source
that Paper~XVI identified are needed for full agreement.

\begin{remark}[QCD angular ordering and the first-jet site]
\label{P19:rem:angular_ordering}
The twelve fraying sites of the trefoil cycle in the sequence
$\cdots\to C\to A\to B\to A\to\cdots$ (with $A$~=~crossing,
$B$~=~midpoint, $C$~=~lobe), with consecutive tangent rotations
$\angle(A,B)=46.75^\circ$ and $\angle(A,C)=83.94^\circ$.
The total tangent rotation across one strand segment
($C\to A\to B\to A$) is $177.44^\circ\approx180^\circ$: the strand
nearly \emph{reverses} the propagation direction, the non-perturbative
topological origin of back-to-back jet kinematics.

QCD angular ordering (coherent parton showers) states that successive
emissions occur at \emph{decreasing} angles from the initial hard
direction.  Applied to the trefoil with threading site at the
A-crossing $t=\pi/6$ ($z=+r_0$), the angular ordering sequence is
determined by the tangent angles of all eleven other sites measured
from the threading tangent $\tang(\pi/6)$:
\begin{equation}
\underbrace{A_{11\pi/6}}_{167.21^\circ}
>\underbrace{C_{2\pi/3}}_{150.22^\circ}
>\underbrace{B_{5\pi/3}}_{141.26^\circ}
>\underbrace{A_{4\pi/3,\,2\pi/3}}_{120.00^\circ}
>\underbrace{B_{\pi},\,C_0}_{{\sim}84^\circ}
>\cdots
\label{P19:eq:angular_ordering_sequence}
\end{equation}
The first jet emerges from the A-crossing at $t=11\pi/6$
($z=-r_0$), the mirror crossing directly below the threading site,
at $\mathbf{167.21^\circ}$ from the threading direction.  This site
is connected to the threading site by the shortest path (arc-length
$\approx L/5$) via a single C-lobe arc going backward.
The $180^\circ-167.21^\circ=\mathbf{12.79^\circ}$ opening angle
(the predicted deviation from perfect back-to-back kinematics) is
set geometrically by the C-lobe curvature $\kappa_C=0.349$.

\textbf{Consistency with the cascade model ordering, and a correction.}
The cascade model of Section~\ref{P19:sec:cascade_falsification}
identifies $E_1\leftrightarrow m_t$ (first exit, most energy) via
forward propagation from the threading site.  Restricting to the
three \emph{upper} A-crossings ($z=+r_0$) that Paper~XVI designates
as the B-departure network, their angles from the threading tangent
are:
\begin{equation}
  \theta(t=\tfrac{\pi}{6}) = 0^\circ,
  \quad
  \theta(t=\tfrac{5\pi}{6}) = 120.00^\circ,
  \quad
  \theta(t=\tfrac{3\pi}{2}) = 120.00^\circ.
  \label{P19:eq:upper_crossing_angles}
\end{equation}
The two non-threading upper crossings are \emph{exactly degenerate}
at $120^\circ$.  This is an exact consequence of Z$_3$ symmetry: the
three upper crossings are separated by $2\pi/3$ in $t$, so the
threading tangent makes exactly $120^\circ$ with both remaining
upper crossings.  The angular ordering confirms the cascade model's
hierarchy --- threading-site region $\to$ heaviest, $120^\circ$
pair $\to$ middle, lower crossings $\to$ lightest --- but
\emph{corrects} it in two ways:
\begin{enumerate}[noitemsep]
\item The cascade model assigns \emph{distinct} energies $E_2\neq E_3$
  to the two intermediate crossings; the angular ordering reveals they
  are \emph{exactly degenerate in angle}, so their mass splitting
  (i.e., the $m_t$--$m_c$ intra-triplet gap) cannot come from the
  forward cascade alone but requires the formation-event Z$_3$-breaking
  azimuthal asymmetry to lift the degeneracy.
\item The cascade assigns the largest energy to the threading site itself;
  the angular ordering identifies the largest-angle emission as the mirror
  crossing ($t=11\pi/6$, $z=-r_0$), connected to threading by the
  backward C-lobe arc.  These are \emph{complementary} descriptions of
  the same top-quark production event: the cascade tracks where energy
  is \emph{deposited}; angular ordering tracks which emission is
  \emph{hardest}.
\end{enumerate}

Among A-crossing sites (up-type quark candidates), the angular
ordering gives the mass hierarchy:
\begin{itemize}[noitemsep]
\item \emph{Heaviest} ($\theta\approx167^\circ$): threading-site region
  ($t=\pi/6$ and its mirror $t=11\pi/6$) --- the top quark.
\item \emph{Middle} ($\theta=120^\circ$, \emph{exactly degenerate}):
  A-crossings at $t=5\pi/6$ and $t=3\pi/2$ (both $z=+r_0$).
  Their exact angular degeneracy is the topological explanation for
  the charm quark: it arises from the Z$_3$-symmetric crossing pair,
  and the small $m_t/m_c$ splitting comes from the formation-event
  azimuthal asymmetry that breaks this degeneracy.
\item \emph{Lightest} ($\theta\leq73^\circ$): lower A-crossings
  ($t=\pi/2$ and $t=7\pi/6$, $z=-r_0$) --- the up quark.
\end{itemize}
Making the angular-ordering hierarchy quantitative --- in particular,
deriving the mass ratios $m_t/m_c$ and $m_c/m_u$ from the curvature
profile $\kappa(t)$ along the formation-event trajectory --- is
an explicit target for the dynamic integrator of Open Problem~O1.
\end{remark}

\begin{remark}[Effect of angular ordering on the reactor angle, quark masses,
  and the $\varepsilon_M$ estimate]
\label{P19:rem:angular_ordering_impact}
Three natural follow-up questions are addressed here.

\textbf{(1) Does the angular ordering change the reactor angle?}
No, directly.  The Paper~XVII estimate
$\sin^2\theta_{13}\sim 3\varepsilon_M\sin^2\theta_M$
(Paper~XVII~\cite{Paper17}, Remark~10.5) %\ref{P17:rem:solar_angle}
uses $\theta_M=31.66^\circ$ (the B-midpoint tangent angle, a fixed
geometric property of the trefoil) and
$\varepsilon_M\approx\kappa_M/(\kappa_A+\kappa_B+\kappa_C)\approx0.022$
(the J$_4$ fraction at midpoints, also fixed by geometry).
Angular ordering tells us \emph{which site emits first}, not the tangent
angles of the sites themselves; both inputs to the reactor-angle formula
are independent of the emission ordering.  At subleading order,
the three B-midpoints are not equally weighted (they appear at
$\theta=141.26^\circ$, $84.56^\circ$, $46.75^\circ$ in the angular-ordering
sequence), so the factor of~$3$ in $3\varepsilon_M\sin^2\theta_M$ is
only a leading-order approximation that the dynamic integrator would
refine.  The C-lobes ($\theta_L=18.8^\circ$ out-of-plane, $|\tang_z^L|=0.321$)
are \emph{not} included in the reactor-angle formula and should not be:
the C-lobes are the hard intra-$\QH=3$ emission channel (high curvature
$\kappa_C=0.349$); they emit into the baryon sector itself, not across to
the $\QH=1$ lepton sector.  Including them would overshoot the observed
$\sin^2\theta_{13}=0.022$ by a factor of~$5$.

\textbf{(2) Does the angular ordering affect the quark mass predictions?}
The two-triplet scale ratio (Result~C,
Proposition~\ref{P19:prop:triplet_ratio}) and the $\sqrt{\vp}$ ladder
(Theorem~\ref{P19:thm:sqrt_phi_ladder}) are unaffected, for different
reasons: the triplet ratio uses the $z$-tangent-squared geometry of
the sites, a fixed property independent of emission ordering; the
$\sqrt{\vp}$ ladder uses only the accepted-as-exact PDG-empirical
inversion and average-scale ratios (Remarks~\ref{P19:rem:cross_ratio},
\ref{P19:rem:tower_structure}), not any site geometry or Jones
polynomial data, so it is likewise unaffected but for a different
reason: it has no ordering-dependent input to begin with.
The angular ordering does refine the physical picture: the exact
$120^\circ$ degeneracy of the two upper A-crossings
(equation~\eqref{P19:eq:upper_crossing_angles}) proves that the
intra-triplet mass splitting between top and charm cannot arise from
an ordering cascade alone but requires the formation-event Z$_3$-breaking
azimuthal asymmetry.  This does not change the predictions but
identifies the dynamical source of the splitting more precisely.

\textbf{(3) Would integrating over the profile function improve $\varepsilon_M$?}
The profile function $f(\rho)=2\arctan(\rho^{-C_2^*})$ satisfies
$f(1)=\pi/2$ \emph{exactly} for any $C_2^*$ (since $\arctan(1^x)=\pi/4$),
so $\sin^4(f(\rho))=1$ at the tube wall $\rho=1$ for all fraying sites.
The transverse profile provides no differential weighting between site
types.  The $\kappa$-profile along each segment arc is the relevant input:
the B-segment curvature drops from $\kappa_A=0.399$ (at bounding
crossings) to $\kappa_B=0.026$ at the midpoint centre and back to
$0.399$, with the arc below $\kappa<0.1$ comprising only
$\approx4.7\%$ of the total trefoil arc.  This is consistent with the
kappa-site proxy $\varepsilon_M\approx0.022$ used in Paper~XVII.
Full profile-function integration would improve $\varepsilon_M$ only
marginally (through 3D overlap corrections at the crossing regions),
not at the leading-order level.
\end{remark}

%══════════════════════════════════════════════════════════════════════════════
\section{Toward the multi-component Jones polynomial framework}
\label{P19:sec:multi_component_jones}
%══════════════════════════════════════════════════════════════════════════════

Paper~XVIII~\cite{Paper18}, Section~4 establishes the topological
reaction rules from single-component Jones polynomials at the WZW
root of unity $q_5=e^{2\pi i/5}$: the rationality-factorisation rule
(Theorem~4.1(ii)) %\ref{P18:thm:jones_rule}
requires that when an irrational-class initial state's reaction
produces a final state containing a rational-class component (here
the $\QH=1$ unknot, $\Jones_U(q_5)=1$), at least one other
final-state component must remain irrational --- satisfied by the
$\QH=3$ trefoil's own irrational satellite invariant
$\Jones_{\rm sat}(q_5)=q_5^3/\vp$.  This enforces the pion emission
$\QH=2+\QH=2\to\QH=3+\QH=1$ as a topological necessity.  The
extension to reactions involving multi-component initial or final
states --- Open Problem~O1 of Paper~XVIII --- requires two
technical inputs: the disjoint-union formula for Jones polynomials,
and the Skein-algebra action on multi-component states.  This
section provides both, sufficient to state the multi-component
reaction rule for the simplest non-trivial channel.

\subsection{The disjoint-union formula}

For an oriented link $L$ that is the disjoint union of two links
$L_1$ and $L_2$ (no crossings between them), the Kauffman bracket
satisfies
\begin{equation}
  \langle L_1 \sqcup L_2 \rangle
  \;=\; (-A^2-A^{-2})\,\langle L_1\rangle\,\langle L_2\rangle,
  \label{P19:eq:kauffman_disjoint}
\end{equation}
where $A=q^{-1/4}$ is the Kauffman variable.  Passing to the Jones
polynomial via $\Jones_L(q)=(-A)^{-3w(L)}\langle L\rangle$ with $w$
the writhe:
\begin{equation}
  \Jones_{L_1\sqcup L_2}(q)
  \;=\; -(q^{1/2}+q^{-1/2})\,\Jones_{L_1}(q)\,\Jones_{L_2}(q),
  \label{P19:eq:jones_disjoint}
\end{equation}
where the factor $-(q^{1/2}+q^{-1/2})=-2\cos(\pi/5)\cdot 2^{1/2}$
at $q=q_5$ equals $-2\cos(2\pi/5)/\cos(\pi/5)\cdot(\ldots)$; a
compact form suffices:
\begin{equation}
  -(q_5^{1/2}+q_5^{-1/2}) \;=\; -\vp.
  \label{P19:eq:disjoint_factor_at_q5}
\end{equation}

\begin{lemma}[Disjoint-union factor at the WZW root]
\label{P19:lem:disjoint_factor}
At $q_5=e^{2\pi i/5}$, the disjoint-union multiplier is $-\vp$.
\end{lemma}

\begin{proof}
$q_5^{1/2}=e^{i\pi/5}$, so
$q_5^{1/2}+q_5^{-1/2}=2\cos(\pi/5)=\vp$.  The overall sign is
$-$ per equation~\eqref{P19:eq:jones_disjoint}.
\qed
\end{proof}

\subsection{Multi-component Jones polynomials for the primary reaction channels}

The primary $\QH=3$ production channels of Paper~XVIII
(Remark~5.4, Remark~5.5) involve multi-component initial or final %\ref{P18:rem:production_conditions}
states:
\begin{align}
  &\text{Two-tube reaction:}
  \quad \QH=2+\QH=2 \;\to\; \QH=3 + \QH=1
  \tag{Ch.~1}\\
  &\text{Satellite reaction:}
  \quad \QH=2+\QH=1 \;\to\; \QH=3
  \tag{Ch.~2}
\end{align}
Both admit spectator configurations in which additional
$\QH=k$ objects propagate through without participating; the
disjoint-union formula supplies the correct Jones-polynomial
counting for such configurations.

\begin{proposition}[Two-component initial state factorisation]
\label{P19:prop:two_comp_initial}
For an initial state consisting of two disjoint $\QH=2$ Hopf
links $T(2,2)_1\sqcup T(2,2)_2$, the Jones polynomial evaluated at
$q_5$ is
\begin{equation}
  \Jones_{T(2,2)\sqcup T(2,2)}(q_5)
  \;=\; -\vp\cdot\Jones_{T(2,2)}(q_5)^2.
  \label{P19:eq:two_hopf_disjoint}
\end{equation}
\end{proposition}

\begin{proof}
Direct application of Lemma~\ref{P19:lem:disjoint_factor} to the
disjoint-union formula.
\qed
\end{proof}

$\Jones_{T(2,2)}(q_5)$ is irrational
(Theorem~\ref{P19:thm:channel_1}):
$\Jones_{T(2,2)}(q_5)=1+\zeta_5^3$, whose modulus is $1/\vp$.
The disjoint-union factor $-\vp$ is also irrational.  The
product $-\vp\cdot\Jones_{T(2,2)}(q_5)^2$ is therefore irrational,
consistent with the Paper~XVIII~\cite{Paper18} arithmetic: the
$\QH=3+\QH=1$ final state has satellite Jones invariant
$|\Jones_{\rm sat}(q_5)|=1/\vp$ (Theorem~2.3 of Paper~XVIII) %\ref{P18:thm:trefoil_satellite_jones}
and $\Jones_U(q_5)=1$; the irrational content of the initial
state is accounted for by the irrational $\Jones_{\rm sat}$, so
no additional irrational spectator is formally required.
The topology-forcing argument for pion emission rests on
charge conservation and component counting
(Theorem~5.3 of Paper~XVIII), %\ref{P18:thm:pion_forced}
not on a rationality factorisation rule.

\subsection{The reaction-rule algebra}

The multi-component Jones polynomial arithmetic of
Proposition~\ref{P19:prop:two_comp_initial} generalises Paper~XVIII's
integer/irrational selection rule to reactions with multi-component
initial or final states.  The correct formulation, however, is
\emph{not} a numerical equality of Jones products; it is an
equivalence up to a unit in the ring of cyclotomic integers.

\begin{definition}[Effective Jones product]
\label{P19:def:effective_jones}
For a multi-component state $\bigsqcup_i L_i$ consisting of $N$
disjoint link-objects, the \emph{effective Jones product} at $q_5$
is
\begin{equation}
  \Jones_{\mathrm{eff}}\bigl(\bigsqcup_i L_i\bigr)
  \;\equiv\;
  (-\vp)^{N-1}\prod_{i=1}^N \Jones_{L_i}(q_5)
  \;\;\in\; \mathbb{Z}[\zeta_5],
  \label{P19:eq:eff_jones_def}
\end{equation}
where $\zeta_5=e^{2\pi i/5}$ and the factors $(-\vp)^{N-1}$ arise
from the disjoint-union formula
(Lemma~\ref{P19:lem:disjoint_factor}).
\end{definition}

\begin{theorem}[Multi-component reaction rule: unit equivalence]
\label{P19:thm:unit_reaction_rule}
A reaction
$\bigsqcup_i L_i^{\mathrm{init}}
\to
\bigsqcup_j L_j^{\mathrm{final}}$
is topologically permitted at the WZW root of unity if and only if
the effective Jones products of the initial and final states differ
by a unit in $\mathbb{Z}[\zeta_5]$:
\begin{equation}
  \frac{\Jones_{\mathrm{eff}}(\mathrm{init})}
       {\Jones_{\mathrm{eff}}(\mathrm{final})}
  \;=\; u
  \;\;\in\; \mathbb{Z}[\zeta_5]^\times.
  \label{P19:eq:unit_rule}
\end{equation}
The units of $\mathbb{Z}[\zeta_5]$ are
$\pm\zeta_5^k\cdot\vp^n$ for $k\in\{0,1,2,3,4\}$ and
$n\in\mathbb{Z}$.
\end{theorem}

\begin{proof}
The ring $\mathbb{Z}[\zeta_5]$ is a principal ideal domain
(class number~$1$ for $p=5$).  The effective Jones product
$\Jones_{\mathrm{eff}}$ is an algebraic integer in
$\mathbb{Z}[\zeta_5]$ by construction: each $\Jones_{L_i}(q_5)$
is an algebraic integer (a polynomial in $q_5$ with integer
coefficients), and $-\vp=\zeta_5^2+\zeta_5^3\in\mathbb{Z}[\zeta_5]$.

A topological reaction conserves the WZW fusion-rule content of the
state.  At the level of Jones polynomials, conservation of
fusion-rule content means that the ideal
$(\Jones_{\mathrm{eff}}(\mathrm{init}))\subseteq
\mathbb{Z}[\zeta_5]$ equals the ideal
$(\Jones_{\mathrm{eff}}(\mathrm{final}))$.  Since
$\mathbb{Z}[\zeta_5]$ is a PID, equality of principal ideals
$(a)=(b)$ holds if and only if $a/b$ is a unit.
\qed
\end{proof}

\begin{remark}[Relation to Paper~XVIII's rationality-class rule]
\label{P19:rem:rationality_vs_unit}
Paper~XVIII~\cite{Paper18}, Theorem~4.1 states the selection rule %\ref{P18:thm:jones_rule}
at the level of rationality classes (rational integers vs.\
irrational algebraic integers).  Theorem~\ref{P19:thm:unit_reaction_rule}
is strictly stronger: unit equivalence implies same rationality
class (since units are never zero, and the product of an irrational
element with a unit remains irrational), but not conversely.  Two
elements in the same rationality class need not differ by a unit
--- they could differ by a non-unit algebraic integer.  The
unit-equivalence formulation is the correct one for a PID.
\end{remark}

\begin{corollary}[Recovery of Paper~XVIII's single-channel rules]
\label{P19:cor:single_channel_recovery}
In the case $N_{\mathrm{init}}=N_{\mathrm{final}}=1$,
$\Jones_{\mathrm{eff}}=\Jones_{L}(q_5)$ on both sides, and
Theorem~\ref{P19:thm:unit_reaction_rule} reduces to
$\Jones_{L_{\mathrm{init}}}(q_5)/\Jones_{L_{\mathrm{final}}}(q_5)\in
\mathbb{Z}[\zeta_5]^\times$, recovering Paper~XVIII's single-channel
picture as the $N=1$ special case of the general rule.  Both
production channels verified in this paper (Theorems~\ref{P19:thm:channel_1},
\ref{P19:thm:channel_2}) have $N_{\mathrm{init}}\geq2$, so neither
is itself an instance of this reduction; no single-component
$T(2,2)\to T(2,3)$ reaction is physically realised in this
framework (both channels require a spectator or companion
component), so this corollary records the algebraic reduction only,
without a worked single-component numeric example.
\end{corollary}

\subsection{Verification of both $\QH=3$ production channels}

We now verify Theorem~\ref{P19:thm:unit_reaction_rule} for both
production channels of Paper~XVIII~\cite{Paper18},
Remark~5.4. %\ref{P18:rem:production_conditions}

\begin{theorem}[Channel~1 verification: two-tube reaction]
\label{P19:thm:channel_1}
The reaction
$T(2,2)\sqcup T(2,2)\to T(2,3)\sqcup U$
satisfies the unit-equivalence rule of
Theorem~\ref{P19:thm:unit_reaction_rule}, with unit factor
$u=\zeta_5+\zeta_5^4=1/\vp$.
\end{theorem}

\begin{proof}
Using Definition~\ref{P19:def:effective_jones} with the Hopf link's
Jones invariant $\Jones_{T(2,2)}(q_5)=1+\zeta_5^3$
(verified directly by the Kauffman-bracket engine of
\texttt{papers/src\_paper18/kauffman\_bracket\_engine.py}: the
standard Hopf-link polynomial $V(t)=-t^{1/2}-t^{5/2}$ at $t=q_5$
equals $1+\zeta_5^3$, of modulus $1/\vp$; see Paper~XVIII~\cite{Paper18}, %\ref{P18:rem:richer_criterion}
the remark on the confinement criterion in Section~2), the
satellite Jones invariant
$\Jones_{\rm sat}(q_5)=q_5^3/\vp=\zeta_5^3/\vp$
(Paper~XVIII~\cite{Paper18}, Theorem~2.3), %\ref{P18:thm:trefoil_satellite_jones}
$\Jones_U(q_5)=1$, $N_{\rm init}=N_{\rm final}=2$:
\begin{align}
  \Jones_{\mathrm{eff}}(\mathrm{init})
  &= (-\vp)^{1}\cdot\Jones_{T(2,2)}(q_5)^2
  \;=\; (-\vp)(1+\zeta_5^3)^2,
  \label{P19:eq:V_init_ch1}\\[4pt]
  \Jones_{\mathrm{eff}}(\mathrm{final})
  &= (-\vp)^{1}\cdot\frac{\zeta_5^3}{\vp}\cdot1
  \;=\; -\zeta_5^3.
  \label{P19:eq:V_final_ch1}
\end{align}
Expanding the square and using $\zeta_5^6=\zeta_5$ (since $\zeta_5^5=1$):
\begin{equation}
  (1+\zeta_5^3)^2 \;=\; 1+2\zeta_5^3+\zeta_5^6 \;=\; 1+\zeta_5+2\zeta_5^3.
  \label{P19:eq:ch1_square}
\end{equation}
Multiplying by $-\vp=-(1+\zeta_5+\zeta_5^4)$ and reducing every power
$\zeta_5^k$ with $k\geq4$ using the cyclotomic relation
$1+\zeta_5+\zeta_5^2+\zeta_5^3+\zeta_5^4=0$ (equivalently
$\zeta_5^4=-(1+\zeta_5+\zeta_5^2+\zeta_5^3)$) gives
\begin{equation}
  (-\vp)(1+\zeta_5+2\zeta_5^3) \;=\; 1+\zeta_5+\zeta_5^3.
  \label{P19:eq:ch1_numerator}
\end{equation}
Dividing by $\Jones_{\mathrm{eff}}(\mathrm{final})=-\zeta_5^3$ is the same
as multiplying by $-\zeta_5^2$ (since $\zeta_5^5=1$, so
$\zeta_5^{-3}=\zeta_5^2$):
\begin{equation}
  \frac{\Jones_{\mathrm{eff}}(\mathrm{init})}
       {\Jones_{\mathrm{eff}}(\mathrm{final})}
  \;=\; -(1+\zeta_5+\zeta_5^3)\zeta_5^2
  \;=\; -(\zeta_5^2+\zeta_5^3+\zeta_5^5)
  \;=\; -(1+\zeta_5^2+\zeta_5^3)
  \;=\; \zeta_5+\zeta_5^4,
  \label{P19:eq:ch1_unit}
\end{equation}
where the last equality is again the cyclotomic relation. This is a
unit in $\mathbb{Z}[\zeta_5]$: it equals $1/\vp$ exactly (since
$\vp=1+\zeta_5+\zeta_5^4$ gives $1/\vp=\vp-1=\zeta_5+\zeta_5^4$), and
the norm $N_{\mathbb{Q}(\zeta_5)/\mathbb{Q}}(\zeta_5+\zeta_5^4)=1$
(computed as the resultant of $\zeta_5+\zeta_5^4$'s minimal polynomial
with the cyclotomic polynomial $\Phi_5$; verified symbolically in
\texttt{papers/src\_paper19/jones\_unit\_equivalence.py}).
Hence $u=\zeta_5+\zeta_5^4=1/\vp\in\mathbb{Z}[\zeta_5]^\times$.
\qed
\end{proof}

\begin{theorem}[Channel~2 verification: satellite reaction]
\label{P19:thm:channel_2}
The reaction
$T(2,2)\sqcup U\to T(2,3)$
satisfies the unit-equivalence rule of
Theorem~\ref{P19:thm:unit_reaction_rule}, with unit factor
$u=-\vp\cdot\zeta_5$.
\end{theorem}

\begin{proof}
$N_{\mathrm{init}}=2$ (Hopf link and unknot, disjoint),
$N_{\mathrm{final}}=1$ (trefoil satellite):
\begin{align}
  \Jones_{\mathrm{eff}}(\mathrm{init})
  &= (-\vp)\cdot\Jones_{T(2,2)}(q_5)\cdot\Jones_U(q_5) \notag\\
  &= (-\vp)(1+\zeta_5^3),
  \label{P19:eq:V_init_ch2}\\[4pt]
  \Jones_{\mathrm{eff}}(\mathrm{final})
  &= (-\vp)^0\cdot\Jones_{\rm sat}(q_5)
  \;=\; \frac{\zeta_5^3}{\vp}.
  \label{P19:eq:V_final_ch2}
\end{align}
Dividing by $\Jones_{\mathrm{eff}}(\mathrm{final})=\zeta_5^3/\vp$ is the
same as multiplying by $\vp\cdot\zeta_5^{-3}=\vp\zeta_5^2$ (since
$\zeta_5^5=1$), so the ratio carries an extra factor of $\vp$ beyond
the one already present in $\Jones_{\mathrm{eff}}(\mathrm{init})$:
\begin{equation}
  \frac{\Jones_{\mathrm{eff}}(\mathrm{init})}
       {\Jones_{\mathrm{eff}}(\mathrm{final})}
  \;=\; -\vp^2(1+\zeta_5^3)\zeta_5^2.
  \label{P19:eq:ch2_setup}
\end{equation}
Using $\vp=1+\zeta_5+\zeta_5^4$ and reducing with the cyclotomic
relation $1+\zeta_5+\zeta_5^2+\zeta_5^3+\zeta_5^4=0$ throughout,
\begin{equation}
  \vp^2 \;=\; (1+\zeta_5+\zeta_5^4)^2 \;=\; 1-\zeta_5^2-\zeta_5^3,
  \label{P19:eq:phi_squared}
\end{equation}
so that
\begin{equation}
  -\vp^2(1+\zeta_5^3) \;=\; -(1-\zeta_5^2-\zeta_5^3)(1+\zeta_5^3)
  \;=\; \zeta_5+\zeta_5^2,
  \label{P19:eq:ch2_numerator}
\end{equation}
and finally, multiplying by $\zeta_5^2$,
\begin{equation}
  \frac{\Jones_{\mathrm{eff}}(\mathrm{init})}
       {\Jones_{\mathrm{eff}}(\mathrm{final})}
  \;=\; (\zeta_5+\zeta_5^2)\zeta_5^2
  \;=\; \zeta_5^3+\zeta_5^4,
  \label{P19:eq:ch2_unit}
\end{equation}
which reduces, via the cyclotomic relation
($\zeta_5^3+\zeta_5^4=-(1+\zeta_5+\zeta_5^2)$), to $-1-\zeta_5-\zeta_5^2$.
This equals $-\vp\cdot\zeta_5$: expanding
$-\vp\zeta_5=-(1+\zeta_5+\zeta_5^4)\zeta_5
=-(\zeta_5+\zeta_5^2+\zeta_5^5)=-(\zeta_5+\zeta_5^2+1)
=-1-\zeta_5-\zeta_5^2$, the same expression.
The unit $u=-\vp\cdot\zeta_5$ is in $\mathbb{Z}[\zeta_5]^\times$:
$-1$, $\vp=1+\zeta_5+\zeta_5^4$, and $\zeta_5$ are each units, so
their product is a unit (norm $N(-\vp\zeta_5)=N(-1)N(\vp)N(\zeta_5)=1$,
verified as the resultant of its minimal polynomial with $\Phi_5$;
verified symbolically in
\texttt{papers/src\_paper19/jones\_unit\_equivalence.py}).
\qed
\end{proof}

\begin{remark}[The unit factors have physical content]
\label{P19:rem:unit_meaning}
The unit factors $u_1=\zeta_5+\zeta_5^4=1/\vp$ (Channel~1, real) and
$u_2=-\vp\cdot\zeta_5$ (Channel~2) carry physical information beyond
the bare ``permitted/forbidden'' classification.  Unlike a
previous draft computation (using the incorrect Hopf-link Jones
value), the two channels do \emph{not} differ by a pure root of
unity: their moduli are reciprocals, $|u_1|=1/\vp$ and $|u_2|=\vp$,
so
\begin{equation}
  \frac{u_2}{u_1} \;=\; -\vp^2\zeta_5,
  \label{P19:eq:channel_ratio}
\end{equation}
an exact identity (verified directly from $u_1,u_2$ above), differing
from a pure phase by the real factor $\vp^2$.  Channel~2
($T(2,2)\sqcup U\to T(2,3)$, $N_{\mathrm{init}}=2\to N_{\mathrm{final}}=1$)
carries one more power of the disjoint-union factor $-\vp$ relative
to its final state than Channel~1
($T(2,2)\sqcup T(2,2)\to T(2,3)\sqcup U$, $N_{\mathrm{init}}=2\to
N_{\mathrm{final}}=2$) does relative to its own --- this asymmetry in
component-count change is the direct algebraic source of the
$\vp^2$ factor (Definition~\ref{P19:def:effective_jones}'s
$(-\vp)^{N-1}$ term). Whether this modulus asymmetry between the
two production channels has further physical meaning (e.g.\ relative
reaction rates or energy scales) is open; it is recorded here as an
exact algebraic fact, not interpreted further.
\end{remark}

The verification of both channels closes Paper~XVIII's Open
Problem~O4 and validates the multi-component framework of
Theorem~\ref{P19:thm:unit_reaction_rule}.  The reaction rule is
not a numerical equality of Jones polynomial products but a
unit-equivalence in the ring of cyclotomic integers --- a
strictly stronger algebraic constraint than the rationality-class
rule of Paper~XVIII, Theorem~4.1. %\ref{P18:thm:jones_rule}

%══════════════════════════════════════════════════════════════════════════════
\section{Experimental Status and Predictions}
\label{P19:sec:experiment}
%══════════════════════════════════════════════════════════════════════════════

We collect the paper's falsifiable predictions and constraints.

\begin{enumerate}[leftmargin=2em,label=\textbf{P\arabic*.}]

\item \textbf{Trefoil minor radius from $\QH=2$ profile.}
The relation $r_0=R_0/C_2^*$ (Theorem~\ref{P19:thm:r0_derivation}) is
an exact identity of the framework.  A test would consist of a
higher-precision numerical solution for the $\QH=2$ profile that
sharpens $C_2^*$ beyond four decimal places and confirms that
$r_0$ at the family boundary of the $\QH=3$ gradient flow tracks
$R_0/C_2^*$ within the profile precision.

\item \textbf{Two-triplet mass-scale ratio.}
The prediction $M_{\mathrm{down}}/M_{\mathrm{up}}\approx 0.189$
(Proposition~\ref{P19:prop:triplet_ratio}) uses only the
Frenet-geometric $z$-tangent values at midpoint and lobe sites.
The observed value is $0.157$, a $20\%$ residual consistent
with Paper~XVI~\cite{Paper16}'s systematic undershoot pattern.
Further test: refined $\tang_z^M$, $\tang_z^L$ from higher-resolution
$\QH=3$ configurations at $N=384$ or beyond should tighten the
prediction toward observation if the missing correction is
sub-leading; otherwise its structural origin lies elsewhere.

\item \textbf{Falsification of single-loss cascade.}
Any future model of the intra-triplet quark mass hierarchy must
avoid the naive single-loss cascade of
Section~\ref{P19:sec:cascade_falsification}
(Proposition~\ref{P19:prop:cascade_falsified}), which is quantitatively
inconsistent with observation by three to six orders of magnitude.
This is not a testable prediction of the framework but a boundary
of the search space, in the style of Paper~XVI's negative results.

\item \textbf{Falsification of uniform $\vp$-tower multiplier.}
Any successful mechanism for the intra-triplet hierarchy must
produce different average exponent spacings for the two isospin
triplets (Remark~\ref{P19:rem:tower_structure}), with the observed
ratio between them approximately $\vp$.  A uniform per-generation
$\vp^{2a}$ multiplier is falsified
(Proposition~\ref{P19:prop:phi_tower_falsified}); the correct
mechanism must be more elaborate.

\item \textbf{Multi-component reaction rules (verified).}
Theorem~\ref{P19:thm:unit_reaction_rule} establishes the
unit-equivalence reaction rule for multi-component states, and
Theorems~\ref{P19:thm:channel_1} and~\ref{P19:thm:channel_2} verify it
for both $\QH=3$ production channels.  A specific consequence:
any multi-body baryon decay must satisfy the unit-equivalence
condition in $\mathbb{Z}[\zeta_5]$; a reaction for which the
ratio of effective Jones products is a non-unit algebraic integer
is topologically forbidden.

\end{enumerate}

%══════════════════════════════════════════════════════════════════════════════
\section{Open Problems}
\label{P19:sec:open}
%══════════════════════════════════════════════════════════════════════════════

The results of this paper close two open problems (the $r_0$
derivation and the multi-component channel verification),
substantially resolve a third (the generation-spacing ratio),
identify a new sharp open problem (O6: $m_s/m_d=n_e=20$)
whose resolution would close the chain from $T_{\mathrm{CMB}}$ to
all six quark masses without any PDG input, explicitly
acknowledge the PMNS $q$-series resummation (O7) deferred
from Paper~XVII~\cite{Paper17} to Paper~XX, and record the
strand-collinearity observation (O8) connecting the trefoil
fraying geometry to QCD jet structure as a programme for
Paper~XX.

\begin{enumerate}[leftmargin=2em,label=\textbf{O\arabic*.}]

\item \textbf{Dynamical partition of $E_{\mathrm{pert}}$ across the
twelve sites.}  The intra-triplet quark mass hierarchy corresponds,
under Conjecture~\ref{P19:conj:six_quarks_reformulated}, to the
six characteristic scales of how the incoming perturbation energy
partitions across the trefoil geometry during the
$\QH=2+\QH=1\to\QH=3$ formation event.  The calculation requires
the dynamic integrator infrastructure of Papers~XVII and~XVIII
extended to simulate the formation event with the velocity field
$\mathbf{v}(\mathbf{x},t)$ recorded at each site; the partition
must reproduce the observed $(m_u,m_c,m_t)$ and $(m_d,m_s,m_b)$
patterns given the site tangent values of
equation~\eqref{P19:eq:tangent_z_values}.  This is Paper~XVIII's
Open Problem~O3 reformulated in the dynamical setting; the static
version of the problem (computing per-site break costs from
gradient-flow relaxation) is shown to be insufficient by
Section~\ref{P19:sec:cascade_falsification}.

\item \textbf{The $\sqrt{\vp}$ quark mass-ratio ladder.}
\emph{Substantially resolved in Section~\ref{P19:sec:phi_tower_falsification}.}
Theorem~\ref{P19:thm:sqrt_phi_ladder} shows that the four
intra-triplet generation gaps form a geometric sequence with ratio
$\sqrt{\vp}$, anchored at the strange-quark ratio $m_s/m_d$
(Remark~\ref{P19:rem:strange_anchor}).  Combined with
Paper~XVI's derived $m_s$ (Corollary~\ref{P19:cor:absolute_masses}),
this predicts $m_b=4216\,\mathrm{MeV}$ at $0.8\%$, $m_t/m_c$ at
$6\%$, and $m_c/m_u$ at $19\%$.  The residuals grow with rung
position and match the always-positive, generation-growing
pattern of Paper~XVI~\cite{Paper16}.
The structural origin of $\sqrt{\vp}$ (the inversion ratio common
to both networks) is open: $\sqrt{|\Jones_{T(2,2)}(q_5)|}=1/\sqrt{\vp}$
does not match the ladder ratio $\sqrt{\vp}$ itself
(Remark~\ref{P19:rem:sqrt_phi_meaning}), so this candidate
explanation is not available.  A first-principles derivation
requires the dynamic integrator to compute
the position-dependent attenuation function $\alpha(s)$
along the trefoil and verify that the inversion ratio $r$ converges
to $\sqrt{\vp}$ independently.

\item \textbf{The subleading $z$-tangent from perturbation memory.}
Proposition~\ref{P19:prop:triplet_ratio}'s leading-order prediction
gives $\langle |\tang_z|^2\rangle_{\mathrm{up}}=0$ at the
crossings (horizontal tangent).  A physical
$\QH=3$ trefoil formed from an actual perturbation event carries
azimuthal asymmetry from the formation direction, which contributes
a subleading $z$-tangent at the crossings that we have not
computed.  This subleading contribution is required for the up-type
triplet to have any mass scale at all; its computation from the
formation-event azimuthal asymmetry is an open problem tied to O1
above.

\item \textbf{Verification of the multi-component reaction rule.}
\emph{Resolved in Section~\ref{P19:sec:multi_component_jones}.}
Both production channels of Paper~XVIII are verified algebraically:
Channel~1 ($T(2,2)\sqcup T(2,2)\to T(2,3)\sqcup U$) has unit
factor $u=\zeta_5+\zeta_5^4=1/\vp$
(Theorem~\ref{P19:thm:channel_1}); Channel~2
($T(2,2)\sqcup U\to T(2,3)$) has unit factor
$u=-\vp\cdot\zeta_5$ (Theorem~\ref{P19:thm:channel_2}).
The reaction rule is not numerical equality of Jones products
but unit-equivalence in $\mathbb{Z}[\zeta_5]$
(Theorem~\ref{P19:thm:unit_reaction_rule}), which is strictly stronger
than the rationality-class rule of Paper~XVIII, Theorem~4.1. %\ref{P18:thm:jones_rule}

\item \textbf{Sharpened production thresholds via multi-component Jones.}
Paper~XVIII's Open Problem~O6 asks for quantitative production
thresholds for $\QH=3$.  Theorem~\ref{P19:thm:unit_reaction_rule}
supplies the multi-component algebra needed to compute the
relative rates of the various $\QH=3$ production channels;
combined with the energy scales
$\Delta E_{2+2}=E(\QH=3)+E(\QH=1)-2E(\QH=2)$ and
$\Delta E_{2+1}=E(\QH=3)-E(\QH=2)-E(\QH=1)$
from Paper~XVII~\cite{Paper17}, the rates and thresholds become
computable.  The Paper~XVIII Open Problem~O6 is thus sharpened
into a well-scoped calculation.

\item \textbf{Derivation of $m_s/m_d = \sqrt{2C_2(1/2)}\cdot e^{8\pi/9} = n_e = 20$.}
Remark~\ref{P19:rem:ms_md_integer} shows that the observed ratio
$m_s/m_d=20$ equals, to $0.047\%$, the E$_8$ Coxeter prediction
multiplied by the spin-Casimir factor $\sqrt{2C_2(j=\tfrac12)}
=\sqrt{3/2}=\sqrt{k/2}$, where the last equality follows from
the structural identity $k=4C_2(\tfrac12)$ exact at $k=3$
(equation~\eqref{P19:eq:k_casimir_identity}).  The factor
$\sqrt{2C_2(j)}$ arises because quarks are spin-$\tfrac12$
fermions while the E$_8$ T-phase formula was derived for bosonic
operators; the same conformal weight $h_\psi=\tfrac12$ that
produces the colour level shift $\delta n^{(C)}$ in Paper~XVI
appears here as a spin-statistics multiplicative correction to
the Coxeter ratio.  A first-principles proof that the fermionic
character of quarks inserts a factor $\sqrt{2C_2(j)}$ into
the E$_8$ T-phase mass ratio would derive $m_s/m_d=n_e=20$
without PDG input and close the chain
$T_{\mathrm{CMB}}+\mathrm{geometry}\to$ all six quark mass ratios.

\item \textbf{PMNS $q$-series resummation and mixing angle corrections
  (deferred to Paper~XX).}
Paper~XVII~\cite{Paper17} derives the leading-order PMNS matrix
as an exact permutation
($\mu\to\nu_3$, $\tau\to\nu_2$, $e\to\nu_1$; Theorem~7.1 %\ref{P17:thm:pmns_permutation}
of Paper~XVII) from the primary branching of
$\mathrm{SU}(2)_3\to\mathrm{SU}(2)_1\otimes M(5,6)$, and
computes the level-$1$ branching coefficients
$b_1^e=0$, $b_1^\mu=-1$, $b_1^\tau=+5$ (Proposition~7.6 %\ref{P17:prop:pmns_b1}
of Paper~XVII).
The present paper reaffirms Paper~XVII's leading-order solar angle
(Proposition~\ref{P19:prop:solar_angle_jones}): $\tan\theta_{12}^{(0)}
=1/\vp$ via the $\mathrm{SU}(2)_3$ S-matrix; the same modulus $1/\vp$
also equals $|\Jones_{T(2,2)}(q_5)|$
(Theorem~\ref{P19:thm:channel_1}), a numerical coincidence recorded
but not established as a derivation (Remark~\ref{P19:rem:solar_physical}).
What remains open is the full set of corrections:
$\Delta\theta_{12}\approx1.69^\circ$ (gap between $31.72^\circ$ and
the observed $33.41^\circ$), the atmospheric deviation
$\Delta\sin^2\theta_{23}\approx0.057$, and the reactor angle
$\sin^2\theta_{13}\approx0.022\approx1/(k+2)^2$.
All three corrections require either RG running from
$\Lambda_{\mathrm{cond}}$ to $m_Z$ (Mechanism~A) or
condensate background mixing from the $\QH=3$ baryon sector
(Mechanism~B) --- neither of which is computable from the present
paper's toolkit.  Paper~XX's responsibility is this calculation.

\item \textbf{Strand structure, QCD angular ordering, and the jet hierarchy.}
The 12-site fraying geometry of the trefoil has a site sequence
$\cdots\to C\to A\to B\to A\to\cdots$ with tangent rotations
$\angle(A,B)=46.75^\circ$ (collinear direction) and
$\angle(A,C)=83.94^\circ$ (hard direction)
(Paper~XVI~\cite{Paper16}; see Remark~\ref{P19:rem:angular_ordering}).
The total tangent rotation across one strand is $177.44^\circ\approx180^\circ$.
The B-midpoints ($\kappa_B\approx0$) provide the collinear propagation
channel; the C-lobes ($\kappa_C\approx0.349$) provide the wide-angle
channel.  These correspond to the collinear and hard gluon emission
channels of perturbative QCD, with the non-zero $\kappa_B=0.026$
acting as a topological infrared cutoff (O8a).

Applying QCD angular ordering (first emission at the largest angle):
given threading at the A-crossing $t=\pi/6$, the first-jet site is
the A-crossing at $t=11\pi/6$ at $167.21^\circ$ --- connected via
the short C-lobe arc backward.  The $12.79^\circ$ opening angle
(deviation from perfect back-to-back kinematics) is geometrically
determined by $\kappa_C$.  Among A-crossings, the angular ordering
gives: heaviest (t-quark, $\theta=167^\circ$), middle-degenerate
(c-quark pair, $\theta=120^\circ$ exactly by Z$_3$), lightest
(u-quark, $\theta\leq73^\circ$).  The $120^\circ$ degeneracy of
the charm pair is a direct prediction of the Z$_3$-symmetric trefoil
(O8b).

Making this identification quantitative ---
(a) deriving $P_{q\to qg}(z)$ from $\kappa(t)$
    and connecting $\kappa_B$ to $\Lambda_{\mathrm{QCD}}$ (O8a), and
(b) computing the energy fractions deposited at each A-crossing
    site from the dynamic integrator and confirming they reproduce
    the $m_t:m_c:m_u$ ratio (O8b)
--- is a programme for Paper~XX.

\end{enumerate}

%══════════════════════════════════════════════════════════════════════════════
\section{Summary}
\label{P19:sec:summary}
%══════════════════════════════════════════════════════════════════════════════

Paper~XVI~\cite{Paper16} established, through systematic negative
results, that no combinatorial or representation-theoretic
mechanism \emph{alone} can supply the generation-dependent
input required for the quark mass hierarchy: the correct
mechanism must supply a genuinely dynamical, non-combinatorial
correction on top of any representation-theoretic backbone.
Paper~XVI's Discussion~I notes that the lepton tower level
$n_e=20$ is fixed by a thermodynamic matching condition
(Paper~XII), not by $\twoT$ group structure, and conjectures an
analogous dynamical origin for the quark sector.  The present
paper both confirms and refines this picture.  On the one hand,
the E$_8$ Coxeter mechanism's prediction $m_s/m_d=e^{8\pi/9}=16.32$
--- which is representation-theoretic (McKay correspondence) ---
is \emph{not} sufficient alone ($18\%$ short of the observed
integer $20$).  On the other hand, the missing factor
$\sqrt{2C_2(j=\tfrac12)}=\sqrt{3/2}$ has a purely dynamical,
non-combinatorial origin: it is the spin-statistics correction
from the quark being a free Weyl fermion, traced to the universal
conformal weight $h_\psi=\tfrac12$ --- the same weight that
Paper~XVI uses to derive the colour level shift $\delta
n^{(C)}=\tfrac12$.  The correct mechanism is therefore
\emph{E$_8$ Coxeter backbone $+$ fermionic correction}:
representation theory supplies the bulk, and quark statistics
supplies the correction.  This is consistent with Pattern~(I) of
Paper~XVI: the dynamical, non-combinatorial piece is essential;
it is just not the \emph{entire} mechanism.

The present paper takes this seriously and reframes the
problem in light of two facts of Paper~XVIII:

\begin{enumerate}[label=(\arabic*)]
\item The $\Zthree$ symmetry of the ideal $T_{2,3}$ trefoil is
broken from the outset by the single-site threading event of the
$\QH=2+\QH=1\to \QH=3$ formation; a physical $\QH=3$ trefoil
carries a memory of its formation location.

\item The $\QH=3$ configuration is transient: gradient-flow searches
(Papers~XV/XVI) find no stable $\QH=3$ saddle point, consistent with
a flat-bottomed energy landscape with no restoring force (Remark~10.2
of Paper~XVIII). %\ref{P18:rem:saddle_difficulty}
It forms, cascades, and decays over the collision timescale, so its
observables are properties of the formation-and-decay transit, not of
any static equilibrium.  (The trefoil's own satellite Jones invariant
is irrational, modulus $1/\vp$ (Theorem~2.3 of Paper~XVIII); %\ref{P18:thm:trefoil_satellite_jones}
rationality is not the source of the transience.)
\end{enumerate}

\begin{remark}[Preliminary numerics beyond $\QH=3$: no distinct equilibrium found]
\label{P19:rem:beyond_qh3_numerics}
Topology-protected constrained gradient flow (20000 steps, angle-clamp
$2^\circ$ per grid point per step, following the method of
Paper~XVI's Q$_H$=3 saddle search) was run on the general $T(2,q_{\rm pol})$
construction at three points spanning the unestablished $\QH=4$--$5$
region beyond the trefoil: $q_{\rm pol}=3.764$ and $4.236$ (the
golden-section division points of $[3,5]$) and $q_{\rm pol}=4.0$ (a
local energy minimum of the unrelaxed construction, coinciding with
the $\QH=4$ sector, which Paper~XVIII acknowledges as charge-conservation-permitted
but establishes no stability or metastability for
(Paper~XVIII~\cite{Paper18}, the remark on higher-charge incoming objects
in the Forced Pion Emission section). %\ref{P18:rem:higher_charge_incoming}
All three runs completed without triggering any dilution or
topology-loss halt condition, but none reached an energy plateau:
all three continued decreasing smoothly through the full step budget,
at comparable rates near the end of the run ($q_{\rm pol}=3.764$ and
$4.236$ reduced by a factor of $\sim12$; $q_{\rm pol}=4.0$ by a factor
of $\sim4.5$, retaining a substantially higher $J_4/K$).  This is
consistent with all three points draining toward the same
topologically trivial vacuum at different rates, rather than any of
them settling near a distinct stable configuration --- the same
flat-bottomed, no-restoring-force picture already established for
$\QH=3$ itself (Remark~10.2 of Paper~XVIII) %\ref{P18:rem:saddle_difficulty}
extending into this unestablished region.  This is a preliminary,
unconverged numerical observation, not a proof of instability at
these $q_{\rm pol}$ values; longer runs or a convergence-based
stopping criterion would be needed to settle the question.
\end{remark}

Quark masses in this framework are therefore not intrinsic energy
scales of a static trefoil but characteristic scales of how a
single incoming perturbation energy partitions across the twelve
trefoil sites during transit.  The paper establishes four
analytic results within this framework:

\begin{enumerate}[label=(\Alph*)]
\item \textbf{The trefoil minor radius is analytically fixed}:
$r_0=R_0/C_2^*$ from the $\QH=2$ profile wall width
(Theorem~\ref{P19:thm:r0_derivation}), matching the numerical
construction value to $0.02\%$.  This closes an open problem
shared by Papers~XV and~XVIII.

\item \textbf{The single-loss cascade is falsified}: neither
$\eta_B=0.81$ (density ratio) nor $\eta_{\mathrm{cross}}=0.266$
(per-crossing transmission) reproduces intra-triplet mass ratios
(Proposition~\ref{P19:prop:cascade_falsified}), miss by three to six
orders of magnitude.

\item \textbf{The uniform $\vp$-tower multiplier is falsified}:
no integer $a$ satisfies the constraint
$m^{(g+1)}/m^{(g)}=\vp^{2a}$ for either triplet
(Proposition~\ref{P19:prop:phi_tower_falsified}).  The correct
replacement structure is a $\sqrt{\vp}$ geometric ladder: the four
intra-triplet generation gaps form a geometric sequence with ratio
$\sqrt{\vp}$, anchored at the strange-quark mass ratio $m_s/m_d$
(Theorem~\ref{P19:thm:sqrt_phi_ladder}).  Combined with
Paper~XVI's derived $m_s=93.32\,\mathrm{MeV}$ ($0.09\%$) and the
framework-derived ratio $m_s/m_d=\sqrt{k/2}\cdot e^{8\pi/9}=19.991$
(equation~\eqref{P19:eq:spin_corrected_ratio}), this yields
$m_b=4214\,\mathrm{MeV}$ ($0.7\%$ off PDG $4183\,\mathrm{MeV}$),
$m_t/m_c$ at $6\%$, and $m_c/m_u$ at $19\%$
(Corollary~\ref{P19:cor:absolute_masses}).  Whether the
Jones-polynomial arithmetic of this paper accounts for the ladder's
common inversion ratio $r=\sqrt{\vp}$ or the up/down generation-spacing
asymmetry is open (Proposition~\ref{P19:prop:jones_modulus_ratio});
the direct comparison does not reproduce the observed ratio.

A key observation (Remark~\ref{P19:rem:ms_md_integer}): the
framework-derived ratio $m_s/m_d=19.991$ is within $0.047\%$ of
the integer $n_e=20$, the lepton CMB tower level the framework
derives independently from $T_{\mathrm{CMB}}$ and geometry.
The predictions use no PDG input; PDG values appear only as
comparison targets.  The dominant residuals (0.7\%, 6\%, 19\%)
come from the ladder structure, not from the base-ratio accuracy.
The spin-Casimir identity $k=4C_2(\tfrac12)$ that makes
$m_s/m_d\approx n_e$ is exact at $k=3$ (Open Problem~O6 seeks a
proof that the 0.047\% gap closes exactly).

\item \textbf{The two-triplet scale ratio is predicted at the
$20\%$ level}: $M_{\mathrm{down}}/M_{\mathrm{up}}\approx 0.189$
from the $z$-tangent-squared network geometry
(Proposition~\ref{P19:prop:triplet_ratio}), matching observation
$0.157$ within the systematic residual identified in
Paper~XVI~\cite{Paper16}.
\end{enumerate}

A fifth result is exact: the multi-component Jones polynomial
framework required for reactions with two or more disjoint
components is developed in
Section~\ref{P19:sec:multi_component_jones}.  The reaction rule is
\emph{unit-equivalence} in the ring
of cyclotomic integers $\mathbb{Z}[\zeta_5]$
(Theorem~\ref{P19:thm:unit_reaction_rule}), not mere rationality-class
equality (Paper~XVIII, Theorem~4.1 %\ref{P18:thm:jones_rule}),
and not numerical equality of Jones products.  Both $\QH=3$ production
channels of Paper~XVIII are verified algebraically:
$T(2,2)\sqcup T(2,2)\to T(2,3)\sqcup U$ with unit factor
$\zeta_5+\zeta_5^4=1/\vp$ (Theorem~\ref{P19:thm:channel_1}), and
$T(2,2)\sqcup U\to T(2,3)$ with unit factor
$-\vp\cdot\zeta_5$ (Theorem~\ref{P19:thm:channel_2}).
This closes Paper~XVIII's Open Problem~O4 and upgrades
Paper~XVIII's rationality-class selection rule to the strictly
stronger unit-equivalence formulation.

A formal result in the lepton sector:
Proposition~\ref{P19:prop:solar_angle_jones} proves that the
leading-order solar mixing angle satisfies
$\tan\theta_{12}^{(0)}=1/\vp=1/d_{1/2}$, where $d_{1/2}$ is the
$\mathrm{SU}(2)_3$ quantum dimension of the spin-$\tfrac12$ primary,
via Paper~XVII's S-matrix.  The same modulus $1/\vp$ also equals
$|\Jones_{T(2,2)}(q_5)|$ (Theorem~\ref{P19:thm:channel_1}), a
numerical coincidence recorded but not established as a derivation
(Remark~\ref{P19:rem:solar_physical}).  Whether $\vp$ also plays a
role in the quark-sector generation hierarchy is open
(Proposition~\ref{P19:prop:jones_modulus_ratio}).

Two conjectures of Paper~XVIII are reconciled: the six-quark
combinatorial (Conjecture~\ref{P19:conj:six_quarks_r19}) and the
writhe-based isospin (Conjecture~\ref{P19:conj:isospin_r19}) are
merged into a single dynamical picture
(Conjecture~\ref{P19:conj:six_quarks_reformulated}) in which the two
isospin triplets correspond to the two networks (crossings vs.\
midpoints/lobes) and the intra-triplet mass hierarchy comes from
the azimuthal position of each site relative to the threading site
of the formation event.

The remaining open problem --- the computation of the dynamical
partition and its comparison to the six quark masses --- is
scoped in Section~\ref{P19:sec:open} and requires the dynamic
integrator infrastructure of Papers~XVII and~XVIII, extended
to simulate the formation event with the velocity field recorded
at each of the twelve sites.

%══════════════════════════════════════════════════════════════════════════════
\begin{thebibliography}{99}

\bibitem{Paper1}
F.~Manfredi,
\textit{The Density-Feedback Faddeev--Niemi Hopfion:
  Exact Algebraic Predictions for Standard-Model Parameters},
Zenodo (2026).
\href{https://doi.org/10.5281/zenodo.19342027}{doi:10.5281/zenodo.19342027}

\bibitem{Paper5}
F.~Manfredi,
\textit{Colour Confinement, the QCD Scale, and Hadronic Structure
  from the Density-Feedback Hopfion},
Zenodo (2026).
\href{https://doi.org/10.5281/zenodo.19638857}{doi:10.5281/zenodo.19638857}

\bibitem{Paper12}
F.~Manfredi,
\textit{The Profile Normalisation Conjecture ---
  Proof via CMB Energy Identity and Two-Route Equivalence},
Zenodo (2026).
\href{https://doi.org/10.5281/zenodo.20210460}{doi:10.5281/zenodo.20210460}

\bibitem{Paper15}
F.~Manfredi,
\textit{The $Q_H=3$ Sector of the Density-Feedback Hopfion},
Zenodo (2026).
\href{https://doi.org/10.5281/zenodo.20691001}{doi:10.5281/zenodo.20691001}

\bibitem{Paper16}
F.~Manfredi,
\textit{Quark Generation Masses in the Density-Feedback Hopfion:
A Survey of Ruled-Out Mechanisms and the $Q_H=3$ Gradient-Flow Landscape},
Zenodo (2026).
\href{https://doi.org/10.5281/zenodo.21012650}{doi:10.5281/zenodo.21012650}

\bibitem{Paper17}
F.~Manfredi,
\textit{The $Q_H=1$ Sector of the Density-Feedback Hopfion:
Topology, Group Structure, Colour Exclusion},
Zenodo (2026).
\href{https://doi.org/10.5281/zenodo.21013412}{doi:10.5281/zenodo.21013412}

\bibitem{Paper18}
F.~Manfredi,
\textit{Preimage Topology and Confinement: The $Q_H=3$ Sector of the Density-Feedback Hopfion},
Zenodo (2026).
\href{https://doi.org/10.5281/zenodo.21047750}{doi:10.5281/zenodo.21047750}

\bibitem{Jones1985}
V.~F.~R.~Jones,
\textit{A polynomial invariant for knots via von Neumann algebras},
Bull.\ Amer.\ Math.\ Soc.\ \textbf{12}, 103--111 (1985).
\href{https://doi.org/10.1090/S0273-0979-1985-15304-2}{doi:10.1090/S0273-0979-1985-15304-2}

\bibitem{Kauffman1987}
L.~H.~Kauffman,
\textit{State models and the Jones polynomial},
Topology \textbf{26}, 395--407 (1987).
\href{https://doi.org/10.1016/0040-9383(87)90009-7}{doi:10.1016/0040-9383(87)90009-7}

\bibitem{Rolfsen1976}
D.~Rolfsen,
\textit{Knots and Links},
Publish or Perish, Berkeley (1976).
ISBN 978-0914098164

\bibitem{PDG2024}
S.~Navas et al.\ (Particle Data Group),
\textit{Review of Particle Physics},
Phys.\ Rev.\ D \textbf{110}, 030001 (2024).
\href{https://doi.org/10.1103/PhysRevD.110.030001}{doi:10.1103/PhysRevD.110.030001}

\end{thebibliography}

\end{document}
